\documentclass[11pt]{amsart}
\usepackage{amsmath,amssymb,amsthm,mathtools,mathrsfs}
\usepackage[margin=1.08in]{geometry}
\usepackage{booktabs}
\usepackage{array}
\usepackage{hyperref}
\hypersetup{hidelinks}
\usepackage[final,expansion=false]{microtype}
\setlength{\parskip}{4pt plus 1pt}
\emergencystretch=2.5em

\newtheorem{theorem}{Theorem}[section]
\newtheorem{proposition}[theorem]{Proposition}
\newtheorem{lemma}[theorem]{Lemma}
\newtheorem{corollary}[theorem]{Corollary}
\theoremstyle{definition}
\newtheorem{definition}[theorem]{Definition}
\newtheorem{problem}[theorem]{Problem}
\theoremstyle{remark}
\newtheorem{remark}[theorem]{Remark}

\newcommand{\C}{\mathbb C}
\newcommand{\R}{\mathbb R}
\newcommand{\N}{\mathbb N}
\newcommand{\dd}{\,d}
\newcommand{\RH}{\mathrm{RH}}
\newcommand{\QW}{\mathcal Q_W}
\newcommand{\HL}{\mathcal H_L}
\newcommand{\K}{\mathcal K}
\newcommand{\XiR}{\Xi}
\newcommand{\eps}{\varepsilon}
\DeclareMathOperator{\Spec}{Spec}
\DeclareMathOperator{\Tr}{Tr}
\DeclareMathOperator{\Var}{Var}

\title[Branch selection and global polarization]
{Branch Selection and Global Polarization in the Connes--CCM Route to RH:\\
Finite Prime--Gamma Inertia, Tate Boundary Geometry,\\
and the Divisor-Sensitivity Obstruction}

\author{David Alejandro Trejo Pizzo}
\thanks{Independent Researcher. \texttt{dtrejopizzo@gmail.com}}

\begin{document}

\begin{abstract}
We organize the Connes--Consani--Moscovici and Suzuki operator programs as
a single route from finite semilocal models to the global Weil criterion.
The first part develops a joint prime--Gamma finite-source refinement.  It
keeps the ordinary von Mangoldt channels, the archimedean channel, the
scalar threshold and the exact radical in one Hermitian system.  Loewner
monotonicity, Haynsworth additivity, scalar Schur innovations, adaptive
Krylov recovery and a heat-core exhaustion remove the finite-selection,
tail and form-core losses.  An endpoint-corrected prolate ladder proves
\(\lvert\mathscr R_L^+\rvert\ll\lambda^7d_4\), while a dyadic divisor
estimate gives
\(\|b_L^+\|\ll\lambda^4\sqrt{(1+L)d_4}\to0\).  The remaining sign in
that realization is the strict physical source surplus over the
anti-shorted deficit.

The second part constructs the local geometric data that this surplus was
trying to express.  A cyclic root Laplacian has pseudodeterminant \(N^2\);
the determinant increment from \(mp^{k-1}\) to \(mp^k\) is \(\log p\),
and its normalized root-stratum incidence is \(p^{-k/2}\).  Their product
is exactly \(\Lambda(p^k)/\sqrt{p^k}\).  At the infinite place, a positive
Gamma spin operator and a two-dimensional polar plane give, by zeta
determinants, the complete factor
\(\frac12s(s-1)\pi^{-s/2}\Gamma(s/2)\).  A shared Tate--Gamma boundary,
matched Abel finite part and covariant charge mixing assemble these data
without discarding prime--prime phases.  All these constructions are
source-defined and use neither zeros nor the Weil sign.

We then determine the precise limitation of the ordinary Hilbert
completion.  The complete charged pushout has absolutely continuous scale
spectrum.  An operator-valued dense-range theorem shows that its Gamma
compliance and unitary charge mixing cannot create point spectrum.
Consequently every equivariant map from the resonant CCM degree one into
this completion annihilates the critical-line eigenclasses.  Thus the
local polarization is valid, but this global completion cannot be the
missing polarization of CCM degree one.

Finally, we synthesize six terminal formulations reached by the route---a
signed limit, a Schur surplus, a Rosati factorization, an equivariant
Hilbert majorant, closure-faithful descent, and a positive global
polarization.  Each contains the complete Weil sign.  Their common
endpoint is analyzed by two complementary obstructions.  Four forgetful
functors are divisor-blind under explicit deformations, while the ordinary
Hilbert pushout fails separately by point-versus-continuous spectral type.
We formulate both mechanisms and isolate two open programs.  One
is geometric: construct an arithmetic intersection product and Hodge-index
polarization on a suitable square of the arithmetic curve.  The other is
analytic but necessarily divisor-sensitive: construct a resonant nuclear
or derived completion retaining evaluation jets before Hilbert reduction.
No proof of either final positivity theorem is claimed.
\end{abstract}

\maketitle
\tableofcontents

\section{Terminal Weil criterion}

Let
\[
 \xi(s)=\frac12s(s-1)\pi^{-s/2}\Gamma(s/2)\zeta(s),
 \qquad
 \XiR(z)=\xi\!\left(\frac12+iz\right).
\]
For a test function \(f\) on \(\R\), let \(W(f)\) denote Weil's explicit
formula distribution, normalized by
\[
\begin{aligned}
 W(f)=&
 \int_\R f(x)(e^{x/2}+e^{-x/2})\,\dd x
 -\sum_{n\ge2}\frac{\Lambda(n)}{\sqrt n}f(\log n)
 -\sum_{n\ge2}\frac{\Lambda(n)}{\sqrt n}f(-\log n)\\
&-(\log 4\pi+\gamma)f(0)
 -\int_0^\infty
 \frac{f(x)+f(-x)-2e^{-x/2}f(0)}{e^{x/2}-e^{-x/2}}\,\dd x .
\end{aligned}
\tag{1.1}
\]
The associated Hermitian quadratic form is
\[
 \QW(v_1,v_2)=W(v_1*\widetilde v_2),
 \qquad \widetilde v(x)=\overline{v(-x)}.
\tag{1.2}
\]

\begin{theorem}[Weil criterion~\cite{Weil1952,Connes1999}]
\label{thm:weil}
The Riemann Hypothesis is equivalent to
\[
 \QW(v,v)\ge0
 \qquad(v\in C_c^\infty(\R)).
\tag{1.3}
\]
\end{theorem}

\begin{remark}
The proof skeleton below is therefore not seeking a new final criterion.
The final criterion is (1.3).  Connes' route is a way of constructing
(1.3) as the limit of finite or localized self-adjoint models.
\end{remark}

\begin{remark}[Attribution and novelty boundary]
The terminal positivity criterion is Weil's criterion, in the
noncommutative-geometric form developed by Connes.  The Hilbert-space
background around Sonin spaces, co-Poisson summation and radical-type
quotients belongs to Burnol, Connes--Consani,
Connes--Consani--Moscovici, Connes--van Suijlekom and Suzuki
\cite{Burnol2004,Burnol2007,CCArch,CCM2025,CvS2025,Suzuki2025,SuzukiScrew}.
Related Connes--Consani developments around the arithmetic curve and
absolute geometry, and Suzuki's Dirichlet-\(L\) variant, are part of the
same operator-theoretic landscape but are not used as hidden inputs below
\cite{CCJacobian,CCAbsolute,SuzukiDirichlet}.
The contribution of this paper starts where these inputs leave an unsigned branch-selection
problem: finite-source Schur inertia, prime-by-prime innovation,
observability staircases, adaptive Krylov recovery, the six-mode prolate
residual estimates, and the heat/hybrid form-core exhaustion.  Classical
linear-algebraic devices used below, such as shorted operators,
Haynsworth inertia additivity, Loewner order and total positivity, are
cited at their first structural use.
\end{remark}

\section{The CCM finite spectral chain}

Put \(L=2\log\lambda\).  The semilocal construction restricts the Weil form
to a finite Fourier space
\[
 V_{L,N}=\operatorname{span}\{e^{2\pi ikx/L}: |k|\le N\}.
\tag{2.1}
\]
Let \(W_{L,N}\) be the corresponding finite Weil matrix.  It contains the
polar, archimedean and prime-power terms of the semilocal trace formula.
The central finite structural facts used by the proof skeleton are those
of the CCM/Connes and Connes--van Suijlekom finite programs
\cite{CCM2025,Connes2026,CvS2025}.

\begin{definition}[Finite CCM package]
\label{def:ccm-package}
For each \(L,N\), the finite package consists of:
\[
 \eps_{L,N}=\min\Spec W_{L,N},
 \qquad
 T_{L,N}=W_{L,N}-\eps_{L,N}I,
\tag{2.2}
\]
the radical line \(\C\xi_{L,N}=\ker T_{L,N}\) when the bottom is simple,
and the positive quotient
\[
 \K_{L,N}=(\C\xi_{L,N})^\perp,
 \qquad
 S_{L,N}=T_{L,N}|_{\K_{L,N}}.
\tag{2.3}
\]
\end{definition}

Connes--Consani--Moscovici identify the finite commutator structure,
construct the quotient operator, and obtain self-adjoint finite spectral
data after shifting and quotienting~\cite{CCM2025}.  Connes' 2026
extremal construction
adds the prolate model \(k_L\), whose transform approximates \(\XiR\).
The component used in the closure theorem is the following finite spectral
statement.

\begin{theorem}[Finite critical-line CCM--Connes model~\cite{CCM2025,Connes2026,CvS2025}]
\label{thm:ccm-input}
For each finite level where the radical quotient is defined, the quotient
operator is self-adjoint in the positive quotient metric.  Its finite
spectral divisor is located on the critical line.  Moreover there is an
explicit prolate model \(k_L\) such that
\[
 \widehat k_L(z)\longrightarrow \XiR(z)
\tag{2.4}
\]
locally uniformly on compact subsets of the open strip
\(|\operatorname{Im}z|<1/2\).
\end{theorem}

\begin{remark}
Theorem~\ref{thm:ccm-input} is the strong finite achievement: the finite
models have real spectral parameters, and the model transform tends to the
right entire function in the interior strip.  What it does not by itself
give is lower-semicontinuity of the indefinite Weil form in the global
limit.
\end{remark}

\section{Localized operator realization}

Suzuki replaces the distributional Weil form by a continuous screw-function
coordinate.  Let \(g\) be the screw function associated with zeta.  On
\((-a,a)\) one obtains a localized quadratic form \({\QW}_a\) and a
self-adjoint operator \(A_a\) associated with it.  The nonlocal
first-order realization provides a concrete operator-theoretic
interpretation of the finite intervals.

\begin{theorem}[Localized self-adjoint Suzuki realization~\cite{Suzuki2025,SuzukiScrew}]
\label{thm:suzuki-input}
For each finite interval \((-a,a)\), the localized Weil quadratic form has
a self-adjoint operator realization \(A_a\).  The construction is
unconditional.  Suzuki's limiting conjecture states that the imaginary
parts of the nontrivial zeros arise from the strong-resolvent limit of
these finite-interval nonlocal realizations as \(a\to\infty\).
\end{theorem}

\begin{remark}
This theorem supplies the correct analytic language for convergence: strong
resolvent convergence of finite-interval operators.  It does not prove the
limit that would put the full zero divisor on the line.
\end{remark}

\section{Finite-to-global branch convergence}

There are two different statements which must not be identified.

\begin{definition}[Divisor convergence]
Let \(\phi_L\) be the actual selected continuum ground branch obtained from
the finite radical quotients after \(N\to\infty\) at fixed \(L\).  Divisor
convergence means that, locally uniformly away from zeros,
\[
 \partial_z^2\log \widehat\phi_L(z)
 \longrightarrow
 \partial_z^2\log \XiR(z).
\tag{4.1}
\]
\end{definition}

\begin{definition}[Positive-form convergence]
Positive-form convergence means that for every \(v\in C_c^\infty(\R)\)
there are finite-level approximants \(v_{L,N}\in V_{L,N}\) with
\[
 v_{L,N}\to v
\tag{4.2}
\]
in a topology strong enough for the explicit formula, and
\[
 \QW(v,v)
 \ge
 \limsup_{L\to\infty}\limsup_{N\to\infty}
 \Bigl\langle T_{L,N}v_{L,N},v_{L,N}\Bigr\rangle
 -
 o(1).
\tag{4.3}
\]
\end{definition}

The finite forms on the right side are nonnegative after shift and
quotient.  Thus positive-form convergence immediately implies Weil
positivity.  Divisor convergence implies the same conclusion if it is
proved with enough control to exclude loss of spectral mass and branch
switching.

\begin{lemma}[Logarithmic-curvature branch form]
\label{lem:branch}
Let \(k_L\) be Connes' prolate model and \(\phi_L\) the actual selected
ground branch.  Since \(\widehat k_L\to\XiR\), divisor convergence is
equivalent, on every simply connected compact set avoiding zeros, to
\[
 \boxed{
 \partial_z^2\log
 \frac{\widehat\phi_L(z)}{\widehat k_L(z)}
 \longrightarrow0 .}
\tag{4.4}
\]
\end{lemma}

\begin{proof}
Choose compatible logarithms on the compact set.  Then
\[
 \partial_z^2\log\widehat\phi_L-\partial_z^2\log\XiR
 =
 \partial_z^2\log\frac{\widehat\phi_L}{\widehat k_L}
 +
 \partial_z^2\log\frac{\widehat k_L}{\XiR}.
\tag{4.5}
\]
The last term tends to zero by local uniform convergence and Cauchy
estimates.  This proves the equivalence.
\end{proof}

\section{Conditional closure of the CCM route}

\begin{theorem}[Connes closure theorem]
\label{thm:closure}
Assume the finite CCM package and Suzuki's localized self-adjoint
realization.  If either positive-form convergence or divisor convergence
with no branch loss holds, then RH follows.
\end{theorem}

\begin{proof}
Under positive-form convergence, the finite shifted quotient forms are
nonnegative:
\[
 \langle T_{L,N}u,u\rangle\ge0
 \qquad(u\perp\xi_{L,N}).
\tag{5.1}
\]
Given \(v\in C_c^\infty(\R)\), choose approximants \(v_{L,N}\) as in
(4.2)--(4.3).  Passing to the limit gives
\[
 \QW(v,v)\ge0.
\tag{5.2}
\]
By Theorem~\ref{thm:weil}, RH follows.

Under divisor convergence with no branch loss, the limiting spectral
divisor of the selected self-adjoint finite operators is the divisor of
\(\XiR\).  Since every finite spectral parameter is real, the limiting
zeros are of the form \(1/2+it\).  Hence all zeros of \(\xi\) lie on the
critical line.  This is RH.
\end{proof}

\begin{remark}
This theorem is the mathematical shape of the Connes proof.  Everything
finite is self-adjoint and critical-line.  The only possible failure is at
the infinite passage: positivity or spectral reality may be lost if the
finite branch does not converge to the correct global branch with sign.
\end{remark}

\section{Localized realization and the global limit}

Suzuki's screw-function operator provides the natural candidate topology for
the limiting passage.  In the notation of Theorem~\ref{thm:suzuki-input},
the desired operator statement is
\[
 A_a \longrightarrow A_\infty
\tag{6.1}
\]
in strong resolvent sense, with
\[
 \Spec(A_\infty)=\{\gamma:\xi(1/2+i\gamma)=0\}
\tag{6.2}
\]
including multiplicities.

\begin{proposition}[Suzuki-to-Connes transfer]
If the CCM quotient operators and the Suzuki finite-interval operators are
identified by a common core and their resolvents differ by \(o(1)\) strongly
on that core, then Suzuki's strong-resolvent limit implies the divisor
convergence needed in Theorem~\ref{thm:closure}.
\end{proposition}

\begin{proof}
Strong resolvent convergence determines spectral projections at continuity
sets of the limiting spectral measure.  The common-core identification
transfers those projections to the CCM branch.  Since the finite CCM
operators are self-adjoint, the limiting spectral measure is real.  If its
spectral determinant is the zeta-regularized determinant of \(\XiR\), then
the divisor of \(\XiR\) is supported on the real spectral line.
\end{proof}

\begin{remark}
This proposition isolates a second version of the same branch-selection step:
one
can prove it as branch convergence of transforms, as form convergence, or
as strong-resolvent convergence plus determinant identification.
\end{remark}

\section{Branch selection as a finite source theorem}

The finite critical-line property and the prolate approximation do not
alone prove RH.  The branch-selection assertion is a signed convergence
theorem.  In
the notation above it can be stated in three equivalent operational forms.

\begin{problem}[Branch selector]
\label{prob:branch}
Let \(\phi_L\) be the actual ground branch selected by the finite CCM
radical quotient after \(N\to\infty\).  Prove that
\[
 \partial_z^2\log
 \frac{\widehat\phi_L(z)}{\widehat k_L(z)}
 \longrightarrow0
\tag{7.1}
\]
locally uniformly on compact subsets of the open critical strip avoiding
zeros.
\end{problem}

\begin{problem}[Positive finite-to-global transfer]
\label{prob:transfer}
For every \(v\in C_c^\infty(\R)\), construct finite quotient approximants
\(v_{L,N}\) for which the completed Weil form satisfies
\[
 \QW(v,v)
 =
 \lim_{L\to\infty}\lim_{N\to\infty}
 \langle T_{L,N}v_{L,N},v_{L,N}\rangle
 +
 R(v),
\tag{7.2}
\]
with \(R(v)\ge0\), or at least with \(R(v)\) controlled so that the right
side is nonnegative.
\end{problem}

\begin{problem}[All-source sharp Doob floor]
\label{prob:doob}
Let \(K>0\) be Riemann's full kernel, normalized by
\(\widehat K=\XiR\), and set
\[
 d\mu_K(x)=2\cosh(x/2)K(x)\,\dd x .
\tag{7.3}
\]
Let \(\mathscr E_K\) be the complete jump energy containing all ordinary
von Mangoldt atoms and the archimedean Gamma jump measure.  Prove, directly
from the literal sources, the sharp inequality
\[
 \boxed{
 \mathscr E_K(r)\ge\frac12\Var_{\mu_K}(r)}
\tag{7.4}
\]
for every even compactly supported smooth \(r\).
\end{problem}

\begin{proposition}[Equivalence of the three gates]
Each of Problems~\ref{prob:branch}, \ref{prob:transfer}, and
\ref{prob:doob}, together with the already established finite CCM--Suzuki
components, implies RH.  Conversely, RH implies each gate in the expected
limiting sense.
\end{proposition}

\begin{proof}
Problem~\ref{prob:branch} gives divisor convergence by
Lemma~\ref{lem:branch}, hence RH by Theorem~\ref{thm:closure}.
Problem~\ref{prob:transfer} gives Weil positivity directly, hence RH by
Theorem~\ref{thm:weil}.  Problem~\ref{prob:doob} is the exact ground-state
form of Weil positivity:
\[
 \QW(Kr,Kr)=\mathscr E_K(r)-\frac12\Var_{\mu_K}(r).
\tag{7.5}
\]
Since every even compactly supported test has the form \(Kr\), this proves
the even Weil criterion, which is equivalent to RH.  Under RH, Weil
positivity holds; the finite critical-line branches then have no off-line
divisor to converge to, and the same identities give the three gates in
their limiting form.
\end{proof}

\section{Finite branch-selection refinement}

The gate which belongs to the Connes route is not the full-kernel
inequality (7.4), because that has already eliminated the finite
semilocal architecture.  The required finite comparison is instead a
one-sided Mosco
transfer for the finite shifted forms, in the sense of Mosco convergence
of closed forms~\cite{Mosco1969,KuwaeShioya2003,Post2012}.  We use Mosco
only as a directional lower-transfer requirement; known instability
phenomena for global path properties under Mosco convergence are a
warning not to infer additional global structure from this step alone
\cite{SuzukiUemura2014}.

Fix \(v\in C_c^\infty(\R)\).  Let \(P_{L,N}v\in V_{L,N}\) be the Fourier
sampling used in the CCM finite model, and remove the finite radical by
\[
 u_{L,N}(v)
 =
 P_{L,N}v-
 \frac{\langle P_{L,N}v,\xi_{L,N}\rangle}
      {\langle \xi_{L,N},\xi_{L,N}\rangle}\xi_{L,N}.
\tag{8.1}
\]
Every object in (8.1) is finite and is defined before any reference to
the zero locations of \(\zeta\).

\begin{definition}[Finite transfer defect]
The finite transfer defect of \(v\) is
\[
 \mathfrak D_{L,N}(v)
 =
 \QW(v,v)-
 \langle T_{L,N}u_{L,N}(v),u_{L,N}(v)\rangle .
\tag{8.2}
\]
Expanding both terms by the explicit formula and by the semilocal trace
formula decomposes it as
\[
 \mathfrak D_{L,N}
 =
 \mathfrak D^{\rm mesh}_{L,N}
 +\mathfrak D^{\rm ext}_{L,N}
 +\mathfrak D^{\Gamma}_{L,N}
 +\mathfrak D^{p}_{L,N}
 +\mathfrak D^{\rm pole}_{L,N}
 +\mathfrak D^{\rm br}_{L,N}.
\tag{8.3}
\]
Here \(\mathfrak D^{\rm mesh}\) is Fourier discretization, 
\(\mathfrak D^{\rm ext}\) is the exterior lattice factor,
\(\mathfrak D^{\Gamma}\), \(\mathfrak D^p\), and
\(\mathfrak D^{\rm pole}\) are the literal Gamma, von Mangoldt and polar
source differences, and \(\mathfrak D^{\rm br}\) is the finite radical
branch correction.
\end{definition}

\begin{proposition}[Known analytic part of the transfer]
\label{prop:known-transfer}
For each fixed \(v\in C_c^\infty(\R)\) there is a cofinal diagonal
\(N=N(L)\), with \(N(L)/L^2\to\infty\), such that
\[
 \mathfrak D^{\rm mesh}_{L,N(L)}(v)
 +\mathfrak D^{\rm ext}_{L,N(L)}(v)
 \longrightarrow0 .
\tag{8.4}
\]
The assertion uses only Fourier approximation, the exterior-lattice
curvature bound, and the fixed-\(L\) cofinal summability already available
in the CCM quotient calculation.
\end{proposition}

\begin{proof}
The Fourier sampling of a fixed compactly supported smooth test converges
rapidly on every fixed physical interval.  The exterior determinant
curvature is \(O_K(L^2/N)\) on compact spectral sets; the assumed diagonal
therefore removes it.  Fixed-\(L\) cofinal summability supplies the
remaining quotient-to-continuum passage before \(L\to\infty\).
\end{proof}

\begin{problem}[Source-side Mosco lower bound]
\label{prob:source-mosco}
Prove that for every \(v\in C_c^\infty(\R)\) one can choose the cofinal
diagonal of Proposition~\ref{prop:known-transfer} so that
\[
 \liminf_{L\to\infty}
 \Bigl(
 \mathfrak D^{\Gamma}_{L,N(L)}(v)
 +\mathfrak D^{p}_{L,N(L)}(v)
 +\mathfrak D^{\rm pole}_{L,N(L)}(v)
 +\mathfrak D^{\rm br}_{L,N(L)}(v)
 \Bigr)
 \ge0 .
\tag{8.5}
\]
\end{problem}

\begin{theorem}[Mosco lower transfer closes the Connes route]
\label{thm:mosco-transfer}
If Problem~\ref{prob:source-mosco} holds, then the positive finite CCM
forms transfer to the completed Weil form.  Consequently RH follows.
\end{theorem}

\begin{proof}
For the diagonal in Proposition~\ref{prop:known-transfer} and
Problem~\ref{prob:source-mosco}, equations (8.2)--(8.5) give
\[
 \mathfrak D_{L,N(L)}(v)\ge-o(1).
\tag{8.6}
\]
Since \(T_{L,N}\ge0\) on the finite quotient, (8.2) implies
\[
 \QW(v,v)
 =
 \langle T_{L,N(L)}u_{L,N(L)}(v),u_{L,N(L)}(v)\rangle
 +\mathfrak D_{L,N(L)}(v)
 \ge-o(1).
\tag{8.7}
\]
Letting \(L\to\infty\) gives \(\QW(v,v)\ge0\).  Since \(v\) was
arbitrary, Weil positivity follows.  Theorem~\ref{thm:weil} then gives
RH.
\end{proof}

\begin{remark}
Problem~\ref{prob:source-mosco} is the next target in the Connes path.  It
is not the classical full-kernel inequality (7.4).  It is a finite-source
statement: the Gamma term, the literal von Mangoldt atoms, the pole and the
radical branch correction must be kept together until after the lower
limit is taken.  Separating them reproduces the previously encountered
sign obstructions.
\end{remark}

\begin{remark}[Joint-source rule]
The source-side lower bound is pursued below through the complete finite
prime--Gamma form and its radical anti-short.  The archimedean, polar,
ordinary-prime and radical contributions may exchange mass, so no
separate lower bound is imposed on any one of them.  The active finite
statement is the joint inertia/gain inequality developed in
Section~\ref{sec:cofinal-inertia}.
\end{remark}

\subsection{Controlled exterior and cross residuals}

We now specify the moving vector used in the closure package and discharge
the even Rayleigh residual.  This specification is essential: a trial
vector satisfying only the two co-Poisson moment conditions may retain a
nonzero endpoint value and hence a literal \(t^{-1}\) exterior carrier.
The following fixed six-mode construction also removes the first two
exterior boundary symbols without changing the \(d_4,d_8\) angle ladder.

Let
\[
 J_6=\{0,4,8,12,16,20\}
\]
and let \(\psi_{j,\lambda}\) be the real normalized even prolate
spheroidal wave functions in the Fourier \(+1\) sector.  Write
\[
 a_j=\psi_{j,\lambda}(0),\qquad
 e_j=\psi_{j,\lambda}(\lambda),\qquad
 q_j=1-\chi_j,\qquad d_j=1-\chi_j^2,
\tag{8.8a}
\]
where \(\chi_j\) is the corresponding finite Fourier eigenvalue, and let
\(\mu_j\) be the separation constant in the prolate differential equation.

\begin{definition}[Endpoint-corrected moving source]
\label{def:endpoint-corrected-source}
Set
\[
 \mathcal K_\lambda^{\rm ep}
 =\left\{
 f=\sum_{j\in J_6}c_j\psi_{j,\lambda}:
 f(0)=0,\quad \int_{\mathbb R}f=0,\quad
 f(\lambda)=f'(\lambda)=0
 \right\}.
\tag{8.8b}
\]
Let \(f_\lambda^{\rm ep}\) be a normalized first generalized eigenvector
of \(\sum d_j|c_j|^2\) on this space.  Apply the co-Poisson construction
and positive-parity projection of Section~8 to \(f_\lambda^{\rm ep}\), and
denote the resulting interior and exterior vectors by
\(V_\lambda^+\) and \(R_\lambda^+\).  Henceforth the moving vector is
\[
 q_L^+=\frac{V_\lambda^+}{\|V_\lambda^+\|}.
\tag{8.8c}
\]
This is also the vector used in the cross residual of
Theorem~\ref{thm:cross-residual-vanishing}.
\end{definition}

\begin{lemma}[The endpoint constraints preserve the angle ladder]
\label{lem:endpoint-angle-ladder}
For all sufficiently large \(\lambda\), the four constraints in
(8.8b) are independent.  If
\(\delta_{0,\lambda}^{\rm ep}\le
\delta_{1,\lambda}^{\rm ep}\) are the generalized eigenvalues of
\(\sum d_j|c_j|^2\) relative to \(\sum|c_j|^2\) on
\(\mathcal K_\lambda^{\rm ep}\), then
\[
 \delta_{0,\lambda}^{\rm ep}\asymp d_4,
 \qquad
 \delta_{1,\lambda}^{\rm ep}\asymp d_8.
\tag{8.8d}
\]
Moreover
\[
 \|V_\lambda^+\|\ge c_V>0
\tag{8.8e}
\]
with \(c_V\) independent of \(\lambda\).
\end{lemma}

\begin{proof}
The exact endpoint-slope identity for a prolate linear combination is
\[
 2\lambda f'(\lambda)
 =\sum_{j\in J_6}c_j(\mu_j-4\pi^2\lambda^4)e_j.
\tag{8.8f}
\]
After imposing \(f(\lambda)=0\), the four constraint rows are therefore
\[
 (a_j),\qquad(q_ja_j),\qquad(e_j),\qquad(\mu_je_j).
\tag{8.8g}
\]
For \(j=4r\), \(0\le r\le5\), the fixed-order prolate-to-Hermite and
endpoint asymptotics give
\[
 |a_{4r}|\asymp1,\qquad
 |e_{4r}|\asymp\lambda^{1/2}q_{4r}^{1/2},\qquad
 \mu_{4r}=\lambda^2(M_0+M_1r+O(\lambda^{-2})),
\tag{8.8h}
\]
where \(M_1\ne0\), together with
\[
 q_{4(r+1)}/q_{4r}\asymp\lambda^8.
\tag{8.8i}
\]
After multiplication by uniformly invertible diagonal row and column
scalings, and with \(y=\lambda^4\), the constraint matrix is a uniform
perturbation of
\[
 \begin{pmatrix}
 1&1&1&1&1&1\\
 1&y^2&y^4&y^6&y^8&y^{10}\\
 1&y&y^2&y^3&y^4&y^5\\
 0&y&2y^2&3y^3&4y^4&5y^5
 \end{pmatrix}.
\tag{8.8j}
\]
Its consecutive four-column minors are nonzero.  Elimination gives two
kernel profiles, up to uniformly bounded nonzero factors,
\[
 \begin{aligned}
 v^{(0)}&=(1,-1,2y^{-1},-y^{-2},y^{-4},0)+O(y^{-2}),\\
 v^{(1)}&=(0,1,-1,2y^{-1},-y^{-2},y^{-4})+O(y^{-2}).
 \end{aligned}
\tag{8.8k}
\]
Since \(d_{4r}\asymp d_4y^{2(r-1)}\), the quotient of the first
profile is \(\asymp d_4\).  After removing that direction, the quotient
of the second is \(\asymp d_8\).  The pivot equations and
Cauchy--Schwarz give the matching lower bounds; the perturbation in
(8.8j) is absorbed by the min--max principle.  This proves (8.8d).
Finally, the first profile has the same nonzero \(0/4\) Hermite limit as
the original moving vector.  Positive-parity co-Poisson projection is
continuous on this fixed-mode limit, which proves (8.8e).
\end{proof}

Put
\[
 \ell_\lambda=(1-P_\lambda)\widehat f_\lambda^{\rm ep}.
\tag{8.8l}
\]

\begin{lemma}[Uniform exterior two-jet estimate]
\label{lem:uniform-exterior-two-jet}
Uniformly for \(t\ge\lambda\),
\[
 \boxed{
 |\ell_\lambda(t)|+
 \lambda^{-1}|\ell_\lambda'(t)|
 \ll \lambda^{7/2}\sqrt{d_4}\,t^{-3}.}
\tag{8.8m}
\]
\end{lemma}

\begin{proof}
Put \(c=2\pi\lambda^2\) and \(x=t/\lambda\).  The Fourier continuation
of each fixed prolate mode is its radial PSWF.  We apply the uniform Bessel
and Liouville--Green expansions of Dunster~\cite{DunsterPSWF}, including
their differentiated error bounds.

For \(1\le x\le2\), the uniform Bessel expansion and
Lemma~\ref{lem:endpoint-angle-ladder} give
\[
 |\ell_\lambda(\lambda x)|+
 \lambda^{-1}|\ell_\lambda'(\lambda x)|
 \ll\lambda^{1/2}\sqrt{d_4},
\tag{8.8n}
\]
which implies (8.8m) because \(x^{-3}\ge1/8\).

For \(2\le x\le c\), specialize Dunster's expansion to fixed order
\(j\in J_6\).  Its small parameter satisfies
\[
 \sigma_j^2=\frac{2j+1}{c}+O(c^{-2}).
\tag{8.8o}
\]
All indices are congruent modulo four, so the term \(-j\pi/2\) has the
same outgoing orientation.  The Liouville phase has the expansion
\[
 \xi(\sigma,x)+J(\sigma)
 =\sqrt{x^2-1}
 +\sigma^2\left(\frac\pi4-\frac12\operatorname {arcsec}x\right)
 +O\!\left(\sigma^4(1+x^{-1})\right).
\tag{8.8p}
\]
Consequently the radial expansion through order \(c^{-1}\) can be written
\[
 \widehat\psi_{j,\lambda}(\lambda x)
 =e_j\bigl(U_{0,\lambda}(x)+\mu_jU_{1,\lambda}(x)\bigr)
 +E_{j,\lambda}(x),
\tag{8.8q}
\]
where the two \(U\)'s are independent of \(j\), and the uniform Olver
remainder, also after one derivative, is
\[
 |E_{j,\lambda}(x)|+
 \lambda^{-1}|\partial_tE_{j,\lambda}(x)|
 \ll |e_j|\bigl(x^{-3}+c^{-2}x^{-1}\bigr).
\tag{8.8r}
\]
Indeed, (8.8p) makes the mode-dependent phase
\((2j+1)/(2x)+O_j(x^{-3}+c^{-1}x^{-1})\), while the radial amplitude is
\(x^{-1}(1+O_j(x^{-2}+c^{-1}))\).  Thus its constant and affine terms in
\(2j+1\) are precisely the two symbols in (8.8q).  They cancel because
\[
 \sum_jc_je_j=0,
 \qquad
 \sum_jc_j\mu_je_j=0
\tag{8.8s}
\]
by the endpoint constraints.  Since
\(\sum_j|c_je_j|\ll\lambda^{1/2}\sqrt{d_4}\) and
\(c^{-2}x^{-1}\le x^{-3}\) for \(x\le c\), summing (8.8q)--(8.8r)
proves (8.8m) throughout this region.

For completeness, the same cancellation controls the range \(x\ge c\)
without extrapolating the middle expansion.  If the outgoing solution is
written
\[
 y_j(t)=e^{2\pi i\lambda t}
 \sum_{m\ge0}a_{m,j}t^{-m-1}+\text{conjugate},
\]
then, with \(A_j=4\pi^2\lambda^4-\mu_j\), its exact recurrence is
\[
 \begin{aligned}
 2i(2\pi\lambda)m a_{m,j}
 ={}&[m(m-1)+A_j]a_{m-1,j}\\
 &+2i(2\pi\lambda)\lambda^2(m-1)a_{m-2,j}\\
 &-\lambda^2(m-2)(m-1)a_{m-3,j}.
 \end{aligned}
\tag{8.8t}
\]
The first two combined coefficients vanish by (8.8s), and
\[
 \sum_jc_ja_{2,j}
 =-\frac1{8(2\pi\lambda)^2}
   \sum_jc_j\mu_j^2a_{0,j}.
\tag{8.8u}
\]
The resulting Volterra majorant, using \(|\mu_j|\ll\lambda^2\), gives
\[
 |\ell_\lambda(t)|+\lambda^{-1}|\ell_\lambda'(t)|
 \ll\lambda^{5/2}\sqrt{d_4}\,t^{-3},
 \qquad x\ge c,
\tag{8.8v}
\]
which is stronger than (8.8m).  The three regions cover
\([\lambda,\infty)\).
\end{proof}

\begin{theorem}[Even Rayleigh residual]
\label{thm:even-rayleigh-residual}
For the endpoint-corrected moving vector of
Definition~\ref{def:endpoint-corrected-source},
\[
 \boxed{|\mathscr R_L^+|\ll\lambda^7d_4.}
\tag{8.8w}
\]
Thus the first residual exponent is \(p_R=7<8\).
\end{theorem}

\begin{proof}
Define, for \(v\ge\lambda\),
\[
 H_\lambda(v)=v^{1/2}\sum_{n\ge1}\ell_\lambda(nv),
 \qquad D=v\partial_v.
\tag{8.8x}
\]
Lemma~\ref{lem:uniform-exterior-two-jet} and absolute summation imply
\[
 |H_\lambda(v)|\ll
 \lambda^{7/2}\sqrt{d_4}\,v^{-5/2},
 \qquad
 |DH_\lambda(v)|\ll
 \lambda^{9/2}\sqrt{d_4}\,v^{-3/2}.
\tag{8.8y}
\]
For \(|\Re w|<1/2\), put
\[
 A_\lambda(w)=\int_\lambda^\infty
 H_\lambda(v)v^{-w}\,\frac{dv}{v}.
\tag{8.8z}
\]
The direct estimate for bounded ordinate and the identity
\[
 wA_\lambda(w)
 =H_\lambda(\lambda)\lambda^{-w}
 +\int_\lambda^\infty
 DH_\lambda(v)v^{-w}\,\frac{dv}{v}
\tag{8.8aa}
\]
for large ordinate yield, uniformly in the strip,
\[
 |A_\lambda(w)|
 \ll\frac{\lambda^{7/2}\sqrt{d_4}}
          {1+|\Im w|}.
\tag{8.8ab}
\]
The exact positive-parity Mellin coordinate is
\[
 b_\lambda(w)=
 \frac{A_\lambda(w)+A_\lambda(-w)}{\sqrt2},
\tag{8.8ac}
\]
and the zero-side explicit formula gives
\[
 QW(R_\lambda^+,R_\lambda^+)
 =\sum_\rho b_\lambda(\rho-\tfrac12)^2.
\tag{8.8ad}
\]
The local zero-count estimate \(N(T+1)-N(T)\ll\log(2+T)\), together
with (8.8ab), makes this series absolutely convergent and gives
\[
 |QW(R_\lambda^+,R_\lambda^+)|
 \ll\lambda^7d_4.
\tag{8.8ae}
\]
Finally, the exact normalized residual identity and (8.8e) give
\[
 |\mathscr R_L^+|
 =\frac{|QW(R_\lambda^+,R_\lambda^+)|}
        {\|V_\lambda^+\|^2}
 \ll\lambda^7d_4,
\]
as claimed.
\end{proof}

\begin{lemma}[Dyadic complex-frequency Bessel bound]
\label{lem:dyadic-complex-bessel}
Let \(I_L=[-L/2,L/2]\), \(\lambda=e^{L/2}\), and let coefficients indexed
by the nontrivial zeta divisor satisfy
\[
 |a_\rho|\le \frac{B}{1+|\gamma|},
 \qquad \rho=\beta+i\gamma .
\tag{8.8af}
\]
Then the symmetric partial sums of
\[
 h(x)=\sum_\rho a_\rho
 e^{-i\gamma x}e^{(\beta-1/2)x}
\tag{8.8ag}
\]
converge in \(L^2(I_L)\), and
\[
 \boxed{\|h\|_{L^2(I_L)}
 \ll B\sqrt{\lambda(1+L)}.}
\tag{8.8ah}
\]
The implied constant is absolute for the Riemann divisor.
\end{lemma}

\begin{proof}
It is enough to treat the positive ordinates, since the negative ordinates
are obtained by conjugation.  Separate a bounded initial block and put
\[
 \mathcal Z_k=\{\rho:2^k\le\gamma<2^{k+1}\},
 \qquad k\ge0.
\tag{8.8ai}
\]
For \(\rho,\rho'\in\mathcal Z_k\), direct integration gives the Gram
entry
\[
 \begin{aligned}
 G_{\rho,\rho'}
 &=\int_{-L/2}^{L/2}
 e^{(\beta+\beta'-1)x}e^{-i(\gamma-\gamma')x}\,dx,\\
 |G_{\rho,\rho'}|
 &\ll\lambda\min\!\left(L,
 \frac1{|\gamma-\gamma'|}\right),
 \end{aligned}
\tag{8.8aj}
\]
where the second expression is read as \(L\) at coincident ordinates.
Partition the ordinate axis into unit intervals.  The unconditional local
zero count
\[
 N(T+1)-N(T)\ll\log(3+T)
\tag{8.8ak}
\]
shows that the interval containing \(\gamma\), together with its two
neighbors, contributes \(O(\lambda L(1+k))\) to a Gram row.  The unit
intervals at distance \(m\ge2\) contribute
\(O(\lambda(1+k)/m)\).  Hence
\[
 \sup_{\rho\in\mathcal Z_k}
 \sum_{\rho'\in\mathcal Z_k}|G_{\rho,\rho'}|
 \ll\lambda(1+L+k)(1+k).
\tag{8.8al}
\]
Schur's test and Riemann--von Mangoldt now give
\[
 \begin{aligned}
 \left\|\sum_{\rho\in\mathcal Z_k}a_\rho
 e^{-i\gamma x}e^{(\beta-1/2)x}\right\|_2^2
 &\ll\lambda(1+L+k)(1+k)
       \sum_{\rho\in\mathcal Z_k}|a_\rho|^2,\\
 \sum_{\rho\in\mathcal Z_k}|a_\rho|^2
 &\ll B^2\frac{1+k}{2^k}.
 \end{aligned}
\tag{8.8am}
\]
After taking square roots, the triangle inequality over the dyadic blocks
yields
\[
 \|h\|_2
 \ll B\sqrt\lambda
 \sum_{k\ge0}
 \frac{(1+k)\sqrt{1+L+k}}{2^{k/2}}
 \ll B\sqrt{\lambda(1+L)}.
\tag{8.8an}
\]
The same summable majorant makes the symmetric partial sums Cauchy in
\(L^2(I_L)\), proving the lemma.
\end{proof}

\begin{theorem}[Cross residual for the direct branch selector]
\label{thm:cross-residual-vanishing}
For the endpoint-corrected moving vector of
Definition~\ref{def:endpoint-corrected-source},
\[
 \boxed{
 \|b_L^+\|
 \ll\lambda^4\sqrt{(1+L)d_4}.}
\tag{8.8ao}
\]
In particular, \(\|b_L^+\|\to0\).
\end{theorem}

\begin{proof}
Put
\[
 e_L^+=\frac{R_\lambda^+}{\|V_\lambda^+\|}.
\tag{8.8ap}
\]
The exact moving-radical identity gives, for every even core vector
\(g\perp q_L^+\),
\[
 \langle b_L^+,g\rangle=-\QW(e_L^+,g).
\tag{8.8aq}
\]
If \(z_\rho=\rho-1/2\), the positive-parity Mellin coordinate is
\[
 \widehat {R_\lambda^+}(z_\rho)
 =b_\lambda(z_\rho),
 \qquad
 b_\lambda(w)=\frac{A_\lambda(w)+A_\lambda(-w)}{\sqrt2}.
\tag{8.8ar}
\]
Equations (8.8ab), (8.8e), and \(0<\Re\rho<1\) therefore give the
complete-divisor representation
\[
 \QW(e_L^+,g)
 =\sum_\rho c_{\rho,\lambda}\widehat g(z_\rho),
 \qquad
 |c_{\rho,\lambda}|
 \ll\frac{B_\lambda}{1+|\gamma|},
 \quad
 B_\lambda=\lambda^{7/2}\sqrt{d_4}.
\tag{8.8as}
\]
Lemma~\ref{lem:dyadic-complex-bessel} shows that this functional has an
\(L^2(I_L)\) Riesz representative \(h_\lambda\) satisfying
\[
 \|h_\lambda\|_2
 \ll B_\lambda\sqrt{\lambda(1+L)}
 =\lambda^4\sqrt{(1+L)d_4}.
\tag{8.8at}
\]
Taking the supremum of (8.8aq) over the even unit ball orthogonal to
\(q_L^+\) proves (8.8ao).  Finally,
\[
 d_4\asymp\lambda^9e^{-4\pi\lambda^2}
\tag{8.8au}
\]
implies
\[
 \lambda^4\sqrt{(1+L)d_4}
 \asymp\lambda^{17/2}\sqrt{1+L}\,e^{-2\pi\lambda^2}
 \longrightarrow0.
\tag{8.8av}
\]
\end{proof}

Theorem~\ref{thm:cross-residual-vanishing} closes the cross residual
required by the direct asymptotic-inertia criterion.  Together with
Theorem~\ref{thm:even-rayleigh-residual}, it gives
\[
 |\mathscr R_L^+|+\|b_L^+\|\longrightarrow0.
\tag{8.8aw}
\]

\section{Cofinal prime--Gamma inertia matching}
\label{sec:cofinal-inertia}

The preceding moving-source analysis proves the two residual estimates.
The individual source components can exchange mass.  We therefore keep
all sources in one finite Hermitian form.

\subsection{Complete and finite joint heads}

On the centered real-even space
\[
 \mathcal H_0=\left\{q\in L^2(\mu_K):\int q\,\dd\mu_K=0\right\},
\tag{9.1}
\]
define the polarized displacement atom
\[
\begin{aligned}
 \mathcal J_u(q,s)
 =\int_{\mathbb R}K(x)K(x-u)\,
 &\overline{q(x)-q(x-u)}\\
 &\times\bigl(s(x)-s(x-u)\bigr)\,\dd x .
\end{aligned}
\tag{9.2}
\]
For \(X\ge2\), put
\[
\begin{aligned}
 \mathcal A_X(q,s)
 ={}&\mathscr E_\Gamma(q,s)
 +\sum_{p^k\le X}\frac{\Lambda(p^k)}{p^{k/2}}
       \mathcal J_{k\log p}(q,s)
 -\frac12\langle q,s\rangle_{\mu_K},\\
 \mathcal T_X(q,s)
 ={}&\sum_{p^k>X}\frac{\Lambda(p^k)}{p^{k/2}}
       \mathcal J_{k\log p}(q,s).
\end{aligned}
\tag{9.3}
\]
The complete form is
\[
 \mathcal A_\infty
 =\mathscr E_K-\frac12\langle\cdot,\cdot\rangle_{\mu_K},
\qquad
 \boxed{\mathcal A_X=\mathcal A_\infty-\mathcal T_X,\quad
        \mathcal T_X\succeq0.}
\tag{9.4}
\]
Thus the finite heads increase in Loewner order toward the complete form.

Let \(V_M\subset\mathcal H_0\) be a nested finite-dimensional mode space
in the complement of the complete radical.  The elementary choice is
generated by centered projections of
\[
 \chi_z(x)=\frac{\cos(zx)}{\cosh(x/2)}
\tag{9.5}
\]
for complete zero orbits \(z\) in a growing spectral rectangle, with all
multiplicity jets retained.  Let \(Z_M:\mathbb C^{d_M}\to V_M\) be the
synthesis map, and write
\[
 \mathbf N_M=Z_M^*Z_M,\qquad
 \mathbf H_{M,X}=Z_M^*\mathcal A_XZ_M .
\tag{9.6}
\]

\begin{theorem}[Cofinal filter-bank approximation]
\label{thm:cofinal-filter-bank}
There are constants \(c>0\) and \(C_M<\infty\) such that, for \(Y\ge X\),
\[
 \boxed{
 0\preceq
 \mathbf N_M^{-1/2}
 (\mathbf H_{M,\infty}-\mathbf H_{M,Y})
 \mathbf N_M^{-1/2}
 \preceq
 \mathbf N_M^{-1/2}
 (\mathbf H_{M,\infty}-\mathbf H_{M,X})
 \mathbf N_M^{-1/2}
 \preceq C_Me^{-cX}I.}
\tag{9.7}
\]
Consequently, for arbitrary \(\varepsilon_M\downarrow0\), the choice
\[
 X(M)\ge\frac1c
 \left(\log C_M+\log\frac1{\varepsilon_M}\right)
\tag{9.8}
\]
makes the omitted literal Euler bank at most \(\varepsilon_M\) in
normalized operator norm.
\end{theorem}

\begin{proof}
The first two inequalities follow from (9.3): each omitted prime-power
atom is a positive Gram form.  For fixed \(M\), every mode in (9.5) has
imaginary part strictly inside the critical strip.  The theta expansion of
\(K(x)K(x-u)\), with \(u=\log p^k\), is therefore bounded on \(V_M\) by a
fixed polynomial times \(\exp(-c p^k)\).  Summation over prime powers
larger than \(X\), followed by conjugation with \(\mathbf N_M^{-1/2}\),
gives the last inequality.  Solving \(C_Me^{-cX}\le\varepsilon_M\)
gives (9.8).
\end{proof}

Define the normalized bank \(S_{M,X}\) by
\[
 S_{M,X}^*S_{M,X}
 =\mathbf N_M^{-1/2}
  \left(\mathscr E_\Gamma+
  \sum_{p^k\le X}\frac{\Lambda(p^k)}{p^{k/2}}
       \mathcal J_{k\log p}\right)_{V_M}
 \mathbf N_M^{-1/2}.
\tag{9.9}
\]
Then the exact signed filter-bank identity is
\[
 \boxed{
 \mathbf G_{M,X}:=
 \mathbf N_M^{-1/2}\mathbf H_{M,X}\mathbf N_M^{-1/2}
 =S_{M,X}^*S_{M,X}-\frac12I.}
\tag{9.10}
\]
No scalar averaging occurs in (9.10): all modes and all retained sources
remain coupled in one matrix.

\subsection{The strict finite frontier}

Put
\[
 \delta_M=
 \lambda_{\min}\!\left(
 \mathbf N_M^{-1/2}\mathbf H_{M,\infty}
 \mathbf N_M^{-1/2}\right),
\qquad
 X_*(M)=\inf\{X:\mathbf H_{M,X}\succeq0\}.
\tag{9.11}
\]

\begin{theorem}[Exact finite-repair criterion]
\label{thm:exact-finite-repair}
For every fixed \(M\),
\[
 \boxed{X_*(M)<\infty\quad\Longleftrightarrow\quad\delta_M>0.}
\tag{9.12}
\]
If \(\delta_M>0\), then
\[
 X_*(M)\le
 \left\lceil\frac1c\log\frac{C_M}{\delta_M}\right\rceil .
\tag{9.13}
\]
If \(\delta_M=0\), every proper finite head has a negative direction and
its least eigenvalue increases to \(0\) without crossing.  If
\(\delta_M<0\), no finite head repairs the completed negative direction.
\end{theorem}

\begin{proof}
If \(\delta_M>0\), Theorem~\ref{thm:cofinal-filter-bank} gives
\[
 \mathbf H_{M,X}\succeq
 (\delta_M-C_Me^{-cX})\mathbf N_M,
\]
which proves (9.13).  Conversely, suppose
\(\mathbf H_{M,X}\succeq0\).  The omitted tail is positive definite on
\(V_M\): if it vanished on a nonzero \(q\), then the displacement atoms
of two omitted distinct primes would make \(q\) periodic with periods
\(\log p,\log r\).  Their ratio is irrational, so continuity makes \(q\)
constant; centeredness then makes \(q=0\), a contradiction.  Hence
\[
 \mathbf H_{M,\infty}
 =\mathbf H_{M,X}
  +(\mathbf H_{M,\infty}-\mathbf H_{M,X})\succ0.
\]
Finite dimensionality gives \(\delta_M>0\).  Evaluating a finite head on a
minimizer of (9.11) proves the last two assertions.
\end{proof}

\begin{remark}[What discreteness proves]
The finite spectra in (9.11) are discrete, but discreteness alone does not
imply (9.12)'s left side.  The scalar sequence \(-e^{-cX}\) is already a
discrete level which increases to \(0\) and never crosses it.  The
arithmetic theorem must therefore prove a strict completed gap or use the
vanishing compensated frontier below.
\end{remark}

\subsection{Radical shorting and the inertia staircase}

The word shorting is used in the Krein--Anderson--Trapp sense:
one eliminates an indefinite or constrained block through the extremal
shorted/compressed operator before reading the remaining Schur complement
\cite{Krein1947,Anderson1971,AndersonTrapp1975}.  The radical itself is
part of the Sonin/co-Poisson and Weil-distribution Hilbert-space
background cited in Remark~1.2; the explicit finite coordinate
\(r_j=K^{(2j)}/K-4^{-j}\) is the coordinate used here for the completed
radical.

Use a distinct symbol for the exact completed radical:
\[
 r_j=\frac{K^{(2j)}}K-4^{-j},
\qquad
 \mathfrak R_J=\operatorname{span}\{r_1,\dots,r_J\}.
\tag{9.14}
\]
Polarization of the radical identity gives
\[
 \mathcal A_\infty(r,f)=0
 \qquad(r\in\overline{\cup_J\mathfrak R_J}).
\tag{9.15}
\]
In bases of \(V_M\) and \(\mathfrak R_J\), write the omitted-tail blocks
as \(\mathbf T^{VV},\mathbf T^{VR},\mathbf T^{RR}\).  The radical block is
strictly positive:
\[
 \mathbf T^{RR}\succ0.
\tag{9.16}
\]
Indeed, vanishing would again give two incommensurable prime periods,
whereas the largest derivative in a nonzero finite combination of the
\(r_j\) has nonconstant theta asymptotics.

The finite-head matrix is therefore
\[
 \mathbf A_{M,J,X}=
 \begin{pmatrix}
  \mathbf H_M-\mathbf T^{VV}&-\mathbf T^{VR}\\
  -\mathbf T^{RV}&-\mathbf T^{RR}
 \end{pmatrix},
\qquad \mathbf H_M=\mathbf H_{M,\infty},
\tag{9.17}
\]
and its Schur complement is
\[
 \mathbf S_{M,J,X}
 =\mathbf H_M-
 \left(
 \mathbf T^{VV}
 -\mathbf T^{VR}(\mathbf T^{RR})^{-1}\mathbf T^{RV}
 \right).
\tag{9.18}
\]

\begin{theorem}[Finite inertia matching~\cite{Haynsworth1968,Gantmacher1959}]
\label{thm:finite-inertia-matching}
One has
\[
 \boxed{
 n_-(\mathbf A_{M,J,X})
 =J+n_-(\mathbf S_{M,J,X}).}
\tag{9.19}
\]
Consequently,
\[
 \boxed{
 n_-(\mathbf A_{M,J,X})=J
 \Longleftrightarrow
 \mathbf T^{VV}
 -\mathbf T^{VR}(\mathbf T^{RR})^{-1}\mathbf T^{RV}
 \preceq\mathbf H_M.}
\tag{9.20}
\]
\end{theorem}

\begin{proof}
Block Gaussian elimination is a congruence from (9.17) to
\(\mathbf S_{M,J,X}\oplus(-\mathbf T^{RR})\).  Equations
(9.16), Sylvester's law of inertia, and (9.18) give (9.19)--(9.20).
\end{proof}

Define the excess negative index and its frontier by
\[
 \kappa(M,J,X)=n_-(\mathbf A_{M,J,X})-J,
\qquad
 X_*(M,J)=\inf\{X:\kappa(M,J,X)=0\}.
\tag{9.21}
\]
Loewner monotonicity and interlacing give
\cite{Loewner1934,Bhatia1997}
\[
\begin{array}{c|c}
\text{parameter increased}&\text{effect on }\kappa\\ \hline
X&\text{nonincreasing}\\
M&\text{nondecreasing}\\
J&\text{nonincreasing}.
\end{array}
\tag{9.22}
\]
Thus \(X_*(M+1,J)\ge X_*(M,J)\), whereas
\(X_*(M,J+1)\le X_*(M,J)\).  This is the exact mode--prime--radical
staircase: a head ending at \(7\) can close one row and fail when the next
mode is admitted.

\subsection{Geometry of a large-prime channel}

The staircase has a literal signal-processing interpretation.  For
\(\chi_z\) in (9.5), uniformly on fixed compact \(z,w\)-blocks,
\[
 \boxed{
 \frac{\log p}{\sqrt p}\mathcal J_{\log p}(\chi_z,\chi_w)
 =\beta_p\,\overline{A_p(z)}A_p(w)+E_p(z,w),}
\tag{9.23}
\]
where
\[
\begin{aligned}
 A_p(z)
 &=2z\sin\!\left(\frac{z\log p}{2}\right)
 +\tanh\!\left(\frac{\log p}{4}\right)
  \cos\!\left(\frac{z\log p}{2}\right),\\
 \beta_p
 &=C_\Xi^2\pi^3(\log p)p^2e^{-2\pi p}
   (1+p^{-1/2})^{-2},
\end{aligned}
\tag{9.24}
\]
and
\[
 |E_p(z,w)|\le C_B\beta_pp^{b-1},
 \qquad b=\max(|\operatorname{Im}z|,|\operatorname{Im}w|)<\frac12.
\tag{9.25}
\]
The proof is Laplace's method at the midpoint \(x=\frac12\log p\):
\(K(x)K(x-\log p)\) has aperture \(p^{-1/2}\), and the first odd Taylor
coefficient of \(\chi_z\) is exactly \(A_p(z)\).  Hence one large prime is
an approximate rank-one derivative sensor.  In its asymptotic regime,
\(d\) newly exposed independent directions require at least \(d\)
independent prime channels.  This proves the moving-bank geometry but not,
by itself, the uniform lower frame bound in (9.20).

\subsection{Finite observability and the exact gain budget}

The distinction between observability and gain can be made exact.  Every
single literal displacement observes every finite elementary mode space.

\begin{theorem}[One-channel finite observability]
\label{thm:one-channel-observability}
For every \(u>0\) and every nonzero \(q\in V_M\),
\[
 \boxed{\mathcal J_u(q,q)>0.}
\tag{9.25a}
\]
Consequently the atom Gram matrix is positive definite on \(V_M\).  Its
normalized least eigenvalue is bounded away from zero uniformly when
\(u\) ranges over a fixed compact subinterval of \((0,\infty)\).
\end{theorem}

\begin{proof}
The real theta kernel is strictly positive.  Therefore
\(\mathcal J_u(q,q)=0\) would imply
\(q(x)=q(x-u)\) for almost every \(x\), and analyticity extends the equality
to the real line.  Every finite linear combination of the modes (9.5),
including confluent multiplicity jets, tends to zero at both ends because
\(\lvert\operatorname {Im}z\rvert<1/2\).  A periodic function with this
decay is zero.  This proves strict positivity.  Continuity in \(u\) and
compactness of the normalized unit sphere give the uniform assertion.
\end{proof}

Write the normalized Gamma deficit and prime bank as
\[
 \mathbf B_M=\frac12I-
 \mathbf N_M^{-1/2}(\mathscr E_\Gamma)_{V_M}\mathbf N_M^{-1/2},
 \qquad
 \mathbf P_{M,X}=\sum_{p^k\le X}\frac{\log p}{p^{k/2}}
 \mathbf N_M^{-1/2}(\mathcal J_{k\log p})_{V_M}\mathbf N_M^{-1/2}.
\tag{9.25b}
\]
Then \(\mathbf G_{M,X}=\mathbf P_{M,X}-\mathbf B_M\), and congruence by
\(\mathbf P_{M,X}^{-1/2}\) gives
\[
 \boxed{
 \mathbf G_{M,X}\succeq0
 \Longleftrightarrow
 \vartheta_{M,X}:=
 \lambda_{\max}(\mathbf P_{M,X}^{-1/2}
 \mathbf B_M\mathbf P_{M,X}^{-1/2})\le1.}
\tag{9.25c}
\]
Whenever a new prime-power atom is added, Theorem
\ref{thm:one-channel-observability} makes \(\vartheta_{M,X}\) strictly
decrease on the positive-deficit sector.  Norm convergence gives
\[
 \boxed{X_*(M)<\infty\quad\Longleftrightarrow\quad
 \vartheta_{M,\infty}<1.}
\tag{9.25d}
\]
Thus aliasing, finite rank and lack of sampling are no longer possible
obstructions.  The strict inequality in (9.25d) is the gain question; it
cannot be replaced by the sum of the individual least atom eigenvalues,
because that replacement discards the cross-prime alignment.

\subsection{Scalar innovations of the completed staircase}

There is an exact one-scalar recursion for the remaining matrix sign.
Choose a basis \(\phi_1,\phi_2,\ldots\) adapted to the nested spaces and
write
\[
 \mathbf H_M=
 \begin{pmatrix}\mathbf H_{M-1}&c_M\\c_M^*&h_M\end{pmatrix}.
\tag{9.25e}
\]
Assuming \(\mathbf H_{M-1}\succ0\), set
\[
 a_M=\mathbf H_{M-1}^{-1}c_M,
 \qquad
 q_M^*=\phi_M-\sum_{j<M}(a_M)_j\phi_j.
\tag{9.25f}
\]

\begin{theorem}[Completed scalar innovation]
\label{thm:scalar-innovation}
The vector \(q_M^*\) is orthogonal to \(V_{M-1}\) in the complete signed
form and
\[
\boxed{
 \begin{aligned}
 \sigma_M
 &:={\mathcal A}_\infty(q_M^*,q_M^*)
 =h_M-c_M^*\mathbf H_{M-1}^{-1}c_M
 =\frac{\det\mathbf H_M}{\det\mathbf H_{M-1}},\\
 \mathbf H_M\succ0
 &\Longleftrightarrow
 \mathbf H_{M-1}\succ0\ \text{and}\ \sigma_M>0.
 \end{aligned}}
\tag{9.25g}
\]
Moreover the scalar retains the complete physical source decomposition
\[
 \boxed{
 \sigma_M=\mathscr E_\Gamma(q_M^*)
 +\sum_{p^k}\frac{\log p}{p^{k/2}}
       \mathcal J_{k\log p}(q_M^*,q_M^*)
 -\frac12\lVert q_M^*\rVert_{\mu_K}^2.}
\tag{9.25h}
\]
\end{theorem}

\begin{proof}
Block multiplication gives
\({\mathcal A}_\infty(v,q_M^*)=0\) for \(v\in V_{M-1}\) and the first
line of (9.25g).  Congruence by triangular block elimination reduces
\(\mathbf H_M\) to \(\mathbf H_{M-1}\oplus[\sigma_M]\); Sylvester's law
and the block determinant formula give the second line.  Substitution in
(9.3) gives (9.25h).
\end{proof}

Before radical anti-shorting, the strict frontier is finite in every row
precisely when all literal innovations (9.25h) are positive.  This scalar
criterion is an exact unshorted precursor: it keeps Gamma, the threshold,
all ordinary von Mangoldt weights and their directional alignment in one
number for each newly admitted mode.  The active shortened target is
identified below by \(\mathfrak g_{M,J,X}>\delta_{M,J,X}\).

The moving-filter interpretation also has an exact update law.  Suppose
the preceding-mode block \(A\) is positive before the atom
\(n=p^k\) is inserted, write \(w_n=\Lambda(n)/\sqrt n\), let
\(U_n=D_n|_{V_{M-1}}\), and put
\[
 D_nq(x)=\sqrt{K(x)K(x-\log n)}
          \{q(x)-q(x-\log n)\}.
\tag{9.25i}
\]
If \(q_{M,X}^*\) is the adaptive finite-head regression residual and
\(r_n=D_nq_{M,X}^*\), direct Schur minimization gives
\[
 \boxed{
 \sigma_{M,+}-\sigma_{M,-}
 =w_n\left\langle r_n,
 (I+w_nU_nA^{-1}U_n^*)^{-1}r_n\right\rangle>0.}
\tag{9.25j}
\]
Indeed, after writing an old-mode correction as \(y\), the new pivot is
the minimum of
\[
 \sigma_{M,-}+\langle y,Ay\rangle
 +w_n\lVert r_n+U_ny\rVert^2;
\]
the normal equation and the Woodbury identity give (9.25j).  Iterating
over the ordered prime powers yields
\[
 \sigma_{M,X_L}=\sigma_{M,X_0}
 +\sum_{X_0<n_\ell\le X_L}\Delta_{M,\ell},
 \qquad \Delta_{M,\ell}>0.
\tag{9.25k}
\]
Thus the row closes at the first index at which this exact cumulative
innovation exceeds the initial deficit.  The series converges to
\(\sigma_M\); it cannot be replaced by a divergent ``infinitely many
channels'' argument, because theta localization makes the large-prime
gain summable.

Several primes can be kept jointly.  For a finite block \(\mathcal B\),
put
\[
 \mathcal D_{\mathcal B}q
 =\bigoplus_{n\in\mathcal B}\sqrt{\frac{\Lambda(n)}{\sqrt n}}D_nq,
 \quad U_{\mathcal B}=\mathcal D_{\mathcal B}|_{V_{M-1}},
 \quad r_{\mathcal B}=\mathcal D_{\mathcal B}q_{M,X}^*.
\tag{9.25l}
\]
The same minimization proves the block update
\[
 \Delta_{\mathcal B}
 =\langle r_{\mathcal B},
 (I+U_{\mathcal B}A^{-1}U_{\mathcal B}^*)^{-1}
 r_{\mathcal B}\rangle
 \ge\frac{\lVert r_{\mathcal B}\rVert^2}
 {1+\lVert U_{\mathcal B}\rVert^2/\lambda_{\min}(A)}.
\tag{9.25m}
\]
Unlike a sum of atomwise floors, (9.25m) retains complementarity between
prime sensors.  It gives a finite interval-certifiable crossing criterion:
partition the unused prime powers into finitely many blocks and verify that
the sum of their lower bounds exceeds the current negative pivot.

\begin{remark}[Why discrete Fermi levels do not force a crossing]
The crossing still requires the physical theta features in (9.25i).  An
exact three-state reversible countermodel shows why.  Let \(G\) be the
path Laplacian on three vertices with both edge weights \(1/8\), and let
\(\Pi\) project onto the centered plane.  Its centered eigenvalues are
\(1/8,3/8\).  Choose positive numbers \(a_n\), indexed by prime powers,
with \(\sum a_n=1/8\), and factor
\(a_n\Pi=(\Lambda(n)/\sqrt n)D_n^*D_n\).  Every channel is then strictly
observable and every finite head has discrete spectrum, while the
completed eigenvalues are
\[
 0,\qquad\frac14,\qquad\frac12.
\tag{9.25n}
\]
The last vector is an exact threshold radical, but the \(1/4\) bound state
never crosses.  Thus reversibility, discreteness, literal positive weights,
strict observability and norm convergence do not imply frontier finiteness;
the missing estimate must use the actual ordinary-prime theta geometry.
\end{remark}

\subsection{Local detection and cumulative adaptive gain}

The strict observability statement can be localized at the midpoint of
one prime-power displacement.  For \(u,\eta>0\), put \(t=u/2\) and
define
\[
\begin{aligned}
 \mathcal J_{u,\eta}(q,s)
 :=\int_{-\eta}^{\eta}&K(t+y)K(t-y)
 \{q(t+y)-q(t-y)\}\\
 &\times\overline{\{s(t+y)-s(t-y)\}}\,\dd y .
\end{aligned}
\tag{9.25o}
\]
Let \(\mathbf C_M^{\rm loc}(u,\eta)\) be its Gram matrix in
\((\phi_1,\ldots,\phi_M)\).

\begin{proposition}[Local Christoffel detection~\cite{Karlin1968,Ismail2005}]
\label{prop:local-christoffel}
For every \(u,\eta>0\),
\[
 \mathbf C_M^{\rm loc}(u,\eta)\succ0,
 \qquad
 \boxed{\kappa_M(u,\eta)
 =\frac{\det\mathbf C_M^{\rm loc}(u,\eta)}
        {\det\mathbf C_{M-1}^{\rm loc}(u,\eta)}>0.}
\tag{9.25p}
\]
Every affine residual \(q=\phi_M+v\), \(v\in V_{M-1}\), satisfies
\(\mathcal J_{u,\eta}(q,q)\ge\kappa_M(u,\eta)\).  Hence an omitted
prime power \(n=p^k\) gives the adaptive lower bound
\[
 \Delta_n\ge\frac{\Lambda(n)}{\sqrt n}
 \kappa_M(\log n,\eta).
\tag{9.25q}
\]
\end{proposition}

\begin{proof}
If the local integral vanishes, strict positivity of \(K\) gives
\(q(t+y)=q(t-y)\) on an interval.  Analytic continuation makes \(q\)
invariant under reflection about \(t\).  Since \(q\) is even, composition
with reflection about zero makes it \(u\)-periodic.  Every vector in the
finite elementary zero-mode and confluent-jet space tends to zero at
\(+\infty\); therefore \(q=0\).  The local Gram is positive definite.
Minimizing its quadratic polynomial over the coefficient of
\(V_{M-1}\) gives the last Schur complement, which is the determinant
ratio in (9.25p).  Restricting the literal atom to the midpoint aperture
and dropping the nonnegative preceding-mode penalty gives (9.25q).
\end{proof}

The atomwise statement has an exact cumulative version.  Partition the
omitted prime powers into consecutive finite blocks
\(\mathcal B_1,\mathcal B_2,\ldots\).  At stage \(\ell\), let
\(A_{\ell-1}\succ0\) be the preceding-mode block, let \(U_\ell\) be the
old-mode feature of \(\mathcal B_\ell\), and let \(\rho_\ell\) be its
adaptive new-mode response.  Set
\[
 B_\ell=U_\ell A_{\ell-1}^{-1}U_\ell^*,\qquad
 R_\ell=\|\rho_\ell\|^2,
 \qquad
 C_\ell=\|A_{\ell-1}^{-1/2}U_\ell^*\rho_\ell\|^2.
\tag{9.25r}
\]

\begin{proposition}[Cumulative adaptive block gain]
\label{prop:cumulative-gain}
If the initial pivot is \(-\delta\), then after \(L\) blocks
\[
\boxed{
 \sigma_L=-\delta+\sum_{\ell=1}^L g_\ell,
 \qquad
 g_\ell=\langle\rho_\ell,(I+B_\ell)^{-1}\rho_\ell\rangle
 \ge\frac{R_\ell^2}{R_\ell+C_\ell}.}
\tag{9.25s}
\]
Consequently a finite prefix crosses if and only if the completed exact
gain satisfies \(\sum_{\ell\ge1}g_\ell>\delta\).  The sufficient
two-moment test is
\[
 \sum_{\ell\ge1}\frac{R_\ell^2}{R_\ell+C_\ell}>\delta.
\tag{9.25t}
\]
\end{proposition}

\begin{proof}
The block Woodbury identity gives
\(\sigma_\ell-\sigma_{\ell-1}=g_\ell\), so the first identity telescopes.
If \(\nu_\ell\) is the spectral measure of \(B_\ell\) at
\(\rho_\ell\), then
\[
 g_\ell=\int\frac{\dd\nu_\ell(s)}{1+s},\quad
 R_\ell=\int\dd\nu_\ell(s),\quad
 R_\ell+C_\ell=\int(1+s)\dd\nu_\ell(s).
\]
Cauchy--Schwarz gives \(R_\ell^2\le
g_\ell(R_\ell+C_\ell)\).  Monotonicity of the partial sums proves the
finite-selection assertion.
\end{proof}

Neither strict local detection nor (9.25s) by itself forces a crossing.
A nonzero completed radical has negative pivot equal to the omitted tail
at every proper cutoff and acquires exactly that amount only in the
limit.  The radical must therefore be removed before its tail response is
compared with the source deficit.

\subsection{Radical-conditioned tail gain}

Let \(\Phi:\C^{M-1}\to V_{M-1}\) and
\(\Psi:\C^J\to\mathfrak R_J\) be coordinate maps.  Realize the omitted
Euler bank by a Hilbert-space feature \(\mathcal V_X\) satisfying
\[
 \|\mathcal V_Xq\|^2=\mathcal T_X(q,q),
 \qquad
 U=\mathcal V_X\Phi,\quad v=\mathcal V_X\phi_M,
 \quad W_J=\mathcal V_X\Psi.
\tag{9.25u}
\]
Strict tail positivity on the finite radical space gives
\(W_J^*W_J\succ0\).  Put
\[
 \Pi_{J,X}=I-W_J(W_J^*W_J)^{-1}W_J^*.
\tag{9.25v}
\]
Thus \(\Pi_{J,X}\) removes precisely the tail features synthesized by a
finite exact-radical correction.

Write the completed mode matrix on
\(V_{M-1}\oplus\operatorname{span}\{\phi_M\}\) as \(H_\infty\).
Maximizing the finite-head form in the radical coordinate gives the exact
anti-short
\[
 \boxed{
 \widehat H_{M,J,X}=H_\infty-
 \begin{pmatrix}U^*\\v^*\end{pmatrix}
 \Pi_{J,X}\begin{pmatrix}U&v\end{pmatrix}
 =\begin{pmatrix}\widehat A&\widehat c\\
                  \widehat c^*&\widehat h\end{pmatrix}.}
\tag{9.25w}
\]
Assume \(\widehat A\succ0\), and define
\[
 \widehat a=\widehat A^{-1}\widehat c,\qquad
 q_{M,J,X}^*=\phi_M-\Phi\widehat a,\qquad
 \widehat\sigma_{M,J,X}
 =\widehat h-\widehat c^*\widehat A^{-1}\widehat c
 =-\delta_{M,J,X}<0.
\tag{9.25x}
\]
The stationary radical coefficient and full saddle residual are
\[
 \zeta_{M,J,X}^*=-(W_J^*W_J)^{-1}W_J^*\mathcal V_Xq_{M,J,X}^*,
 \qquad
 \widetilde q_{M,J,X}^*=q_{M,J,X}^*+\Psi\zeta_{M,J,X}^*.
\tag{9.25y}
\]
In particular,
\[
 \mathcal V_X\widetilde q_{M,J,X}^*
 =\Pi_{J,X}\mathcal V_Xq_{M,J,X}^*.
\tag{9.25z}
\]

\begin{theorem}[Exact radical-conditioned adaptive gain]
\label{thm:radical-conditioned-gain}
Set
\[
 z_{M,J,X}=\Pi_{J,X}\mathcal V_Xq_{M,J,X}^*,
 \qquad \overline U=\Pi_{J,X}U.
\]
Then the completed last pivot is
\[
 \boxed{
 \sigma_{M,\infty}=-\delta_{M,J,X}+\mathfrak g_{M,J,X},
 \qquad
 \mathfrak g_{M,J,X}
 =\left\langle z_{M,J,X},
 (I+\overline U\widehat A^{-1}\overline U^*)^{-1}
 z_{M,J,X}\right\rangle.}
\tag{9.25aa}
\]
Equivalently,
\[
 \boxed{
 \mathfrak g_{M,J,X}
 =\sup_{\omega\ne0}
 \frac{|\langle\omega,z_{M,J,X}\rangle|^2}
 {\|\omega\|^2+
  \|\widehat A^{-1/2}\overline U^*\omega\|^2}.}
\tag{9.25ab}
\]
Finite-support tail combiners are sufficient in the strict inequality.
Thus
\[
 \boxed{\sigma_{M,\infty}>0
 \quad\Longleftrightarrow\quad
 \mathfrak g_{M,J,X}>\delta_{M,J,X}.}
\tag{9.25ac}
\]
\end{theorem}

\begin{proof}
The radical block in the joint finite-head matrix is
\(-W_J^*W_J\).  Its Schur complement is (9.25w), while its normal
equation gives (9.25y)--(9.25z).  Adding the omitted bank back to
\(\widehat H_{M,J,X}\) is therefore the positive update with old feature
\(\overline U\) and adaptive response \(z_{M,J,X}\).  The block Woodbury
identity gives (9.25aa).  The Rayleigh dual formula for the positive
operator \(I+\overline U\widehat A^{-1}\overline U^*\) gives (9.25ab).
Finite-support vectors are dense in the tail direct sum, and the two
quadratic functionals in (9.25ab) are continuous; a strict witness may
therefore be approximated by a finite-support witness.  Finally
(9.25ac) follows from (9.25aa).
\end{proof}

The formula also has the joint variational form
\[
 \boxed{
 \mathfrak g_{M,J,X}
 =\inf_{\substack{d\in V_{M-1}\\r\in\mathfrak R_J}}
 \left\{\widehat{\mathcal A}(d,d)
 +\|\mathcal V_X(q_{M,J,X}^*+d+r)\|^2\right\},}
\tag{9.25ad}
\]
where \(\widehat{\mathcal A}\) is the positive old-mode form with matrix
\(\widehat A\).  This identity keeps one common radical correction and
one common old-mode correction across all prime powers.

\subsection{Radical-conditioned Christoffel and finite determinant
exhaustion}

On the elementary analytic space \(V_M\oplus\mathfrak R_J\), every
literal midpoint form (9.25o) is positive definite.  Indeed, a null
vector is periodic by the reflection argument.  Its zero-mode component
decays, whereas a nonzero largest radical derivative has a nonconstant
leading polynomial in \(e^{2x}\); their sum cannot be periodic.  Hence,
if \(\mathbf C_{M,J}^{\rm loc}(u,\eta)\) is the local Gram in a basis of
\(V_M\oplus\mathfrak R_J\),
\[
 \boxed{
 \kappa_{M,J}(u,\eta)
 =\frac{\det\mathbf C_{M,J}^{\rm loc}(u,\eta)}
        {\det\mathbf C_{M-1,J}^{\rm loc}(u,\eta)}>0.}
\tag{9.25ae}
\]
Consequently every finite omitted block \(\mathcal B\) gives the valid
but not necessarily sharp certificate
\[
 \mathfrak g_{M,J,X}\ge
 \sum_{n\in\mathcal B}\frac{\Lambda(n)}{\sqrt n}
 \kappa_{M,J}(\log n,\eta_n).
\tag{9.25af}
\]
The separate minima in (9.25af) can lose the joint alignment.  The
following augmented determinant removes that loss.

For \(Y>X\), define
\[
\begin{aligned}
 \mathfrak C_{M,J;X,Y}:=
 \inf_{\substack{d\in V_{M-1}\\r\in\mathfrak R_J}}
 \bigg\{&\widehat{\mathcal A}(d,d)\\
 &+\sum_{\substack{n=p^k\\X<n\le Y}}
 \frac{\Lambda(n)}{\sqrt n}
 \mathcal J_{\log n}(q_{M,J,X}^*+d+r,q_{M,J,X}^*+d+r)
 \bigg\}.
\end{aligned}
\tag{9.25ag}
\]
In the ordered affine basis
\((V_{M-1},\mathfrak R_J;q_{M,J,X}^*)\), let
\(\mathbb G_n\) be the Gram matrix of \(\mathcal J_{\log n}\), and put
\[
 \mathbb C_{M,J;X,Y}
 =\operatorname{diag}(\widehat A,0_J,0)
 +\sum_{\substack{n=p^k\\X<n\le Y}}
 \frac{\Lambda(n)}{\sqrt n}\mathbb G_n.
\tag{9.25ah}
\]
Let \(\mathbb C_{M,J;X,Y}^{-}\) be the principal block obtained by
deleting the last row and column.

\begin{theorem}[Augmented finite determinant exhaustion]
\label{thm:augmented-exhaustion}
For every \(Y\) containing at least one omitted prime power,
\[
 \boxed{
 \mathfrak C_{M,J;X,Y}
 =\frac{\det\mathbb C_{M,J;X,Y}}
        {\det\mathbb C_{M,J;X,Y}^{-}}>0.}
\tag{9.25ai}
\]
Moreover,
\[
 \boxed{
 0<\mathfrak C_{M,J;X,Y}<\mathfrak g_{M,J,X},\qquad
 \mathfrak C_{M,J;X,Y}\nearrow\mathfrak g_{M,J,X},}
\tag{9.25aj}
\]
and, on each fixed finite row, there are \(C,c>0\) and a polynomial
\(Q\) such that
\[
 0\le\mathfrak g_{M,J,X}-\mathfrak C_{M,J;X,Y}
 \le C Q(Y)e^{-cY}.
\tag{9.25ak}
\]
Consequently
\[
 \boxed{
 \sigma_{M,\infty}>0
 \quad\Longleftrightarrow\quad
 \mathfrak C_{M,J;X,Y}>\delta_{M,J,X}
 \text{ for some finite }Y.}
\tag{9.25al}
\]
\end{theorem}

\begin{proof}
A single displacement Gram is positive definite on
\(V_M\oplus\mathfrak R_J\) by the periodicity/asymptotic argument above.
The positive ordinary weights therefore make both the matrix in
(9.25ah) and its preceding principal block positive definite.  Its last
Schur complement is exactly the minimization (9.25ag), proving
(9.25ai).

The objectives in (9.25ag) increase pointwise with \(Y\) and are bounded
above by (9.25ad).  One fixed nonempty head is coercive on the
finite-dimensional correction space, so all finite minimizers lie in a
common compact set.  Passing to a convergent subsequence and then using
monotone convergence of the nonnegative prime-power series proves that
the limit is (9.25ad).  If \((d_\infty,r_\infty)\) minimizes the complete
objective, then \(q_{M,J,X}^*+d_\infty+r_\infty\ne0\); every omitted
literal atom is strictly positive on it.  Every proper cutoff therefore
omits a positive term, proving the strict inequality in (9.25aj).
The fixed-space theta-overlap estimate bounds the remaining Gram tail by
\(CQ(Y)e^{-cY}\); compactness of the minimizers gives (9.25ak).
Finally combine (9.25aa) and (9.25aj): a positive limiting margin is
eventually larger than the vanishing tail loss, while a finite crossing
implies \(\mathfrak g_{M,J,X}>\delta_{M,J,X}\).  This proves (9.25al).
\end{proof}

The finite source equation eliminates the abstract deficit.  Namely,
\[
\boxed{
\begin{aligned}
 \delta_{M,J,X}={}&\frac12\|\widetilde q_{M,J,X}^*\|_{\mu_K}^2
 -\mathscr E_\Gamma(\widetilde q_{M,J,X}^*)\\
 &-\sum_{p^k\le X}\frac{\log p}{p^{k/2}}
 \mathcal J_{k\log p}(\widetilde q_{M,J,X}^*,
                       \widetilde q_{M,J,X}^*).
\end{aligned}}
\tag{9.25am}
\]
There is an exact cancellation which rules out a raw-response shortcut.
Put
\[
 R_{M,J,X}=\|z_{M,J,X}\|^2,
 \qquad
 \mathfrak l_{M,J,X}=R_{M,J,X}-\mathfrak g_{M,J,X}\ge0.
\]
Radical invariance of the complete form, (9.25z), and
\(\mathcal A_X=\mathcal A_\infty-\mathcal T_X\) give
\[
 \boxed{
 R_{M,J,X}-\delta_{M,J,X}
 =\mathcal A_\infty(q_{M,J,X}^*,q_{M,J,X}^*),
 \qquad
 \mathfrak g_{M,J,X}-\delta_{M,J,X}
 =\mathcal A_\infty(q_{M,J,X}^*,q_{M,J,X}^*)
  -\mathfrak l_{M,J,X}.}
\tag{9.25an}
\]
Thus even the prerequisite \(R_{M,J,X}>\delta_{M,J,X}\) is already
positivity of the completed form on the current residual.  Bounds which
first dominate the raw response and then pay a positive regression penalty
are strictly stronger than the required statement.  The augmented
determinant is sharp because it retains \(\mathfrak l_{M,J,X}\) exactly.

The same assertion admits a source-balanced bordered-determinant form.
Write
\[
 \tau_{d+1}(Y)=\det\mathbb C_{M,J;X,Y},\qquad
 \tau_d(Y)=\det\mathbb C_{M,J;X,Y}^{-},
\]
let \(e_*\) denote the last coordinate vector, and set
\[
 \mathbb P_{M,J;X,Y}
 =\mathbb C_{M,J;X,Y}
  -\delta_{M,J,X}e_*e_*^*.
\]
Taking the Schur complement of the preceding principal block gives the
exact identity
\[
 \boxed{
 \det\mathbb P_{M,J;X,Y}
 =\tau_d(Y)\{
   \mathfrak C_{M,J;X,Y}-\delta_{M,J,X}\}
 =\tau_{d+1}(Y)-\delta_{M,J,X}\tau_d(Y).}
\]
Since \(\tau_d(Y)>0\), finite crossing is therefore precisely the strict
minor domination
\[
 \boxed{\tau_{d+1}(Y)>\delta_{M,J,X}\tau_d(Y).}
\]
This recoding introduces neither a localization loss nor a separate
effective-resistance estimate.

Thus (9.25al) is the completely finite inequality
\[
\begin{aligned}
 \frac{\det\mathbb C_{M,J;X,Y}}
      {\det\mathbb C_{M,J;X,Y}^{-}}
 &+\mathscr E_\Gamma(\widetilde q_{M,J,X}^*)\\
 &+\sum_{p^k\le X}\frac{\log p}{p^{k/2}}
 \mathcal J_{k\log p}(\widetilde q_{M,J,X}^*,
                       \widetilde q_{M,J,X}^*)
 >\frac12\|\widetilde q_{M,J,X}^*\|_{\mu_K}^2.
\end{aligned}
\tag{9.25ao}
\]
It retains Gamma, the scalar threshold, the exact radical correction, the
old-mode adaptation, every retained ordinary prime power, and one finite
omitted block in a single determinant.

\subsection{Signed amplitude transfer and adaptive Krylov recovery}
\label{sec:signed-krylov}

Fix one row \((M,J,X,Y)\), put \(d=M-1+J\), and define the individual
prime-power feature maps by
\[
 (\mathcal D_nq)(x)=
 \left(\frac{\Lambda(n)}{\sqrt n}K(x)K(x-\log n)\right)^{1/2}
 \bigl(q(x)-q(x-\log n)\bigr),
 \qquad n=p^k.
\]
Set
\[
 \mathcal D_{X,Y}=\bigoplus_{\substack{n=p^k\\X<n\le Y}}\mathcal D_n.
\]
This is a square-root feature map for the finite omitted prime-power
block:
\[
 \|\mathcal D_{X,Y}q\|^2
 =\sum_{\substack{n=p^k\\X<n\le Y}}
 \frac{\Lambda(n)}{\sqrt n}\mathcal J_{\log n}(q,q).
\tag{9.25ap}
\]
Realize its observation space as \(L^2(\Omega_Y,\nu_Y)\), where
\(\Omega_Y\) is the disjoint union of the channel copies of \(\mathbb R\).
The one-particle exterior space is
\[
 \mathcal H_Y=\mathbb C^{M-1}\oplus L^2(\Omega_Y,\nu_Y)
 =L^2(\widetilde\Omega_Y,\widetilde\nu_Y),
\]
where \(\widetilde\Omega_Y\) adjoins \(M-1\) counting-measure atoms.
We use the normalized exterior convention
\[
 \|\Phi\|_{[r]}^2=\frac1{r!}
 \int_{\widetilde\Omega_Y^r}|\Phi|^2\,\dd\widetilde\nu_Y^r.
\]
Abbreviate
\[
 A=\widehat A,\qquad
 U=\mathcal D_{X,Y}\Phi,\qquad
 W=\mathcal D_{X,Y}\Psi,\qquad
 v=\mathcal D_{X,Y}q_{M,J,X}^* .
\]
The radical basis and finite observability give \(W^*W\succ0\).  Set
\[
 \Pi_Y=I-W(W^*W)^{-1}W^*,\qquad
 z=\Pi_Yv,\qquad
 B=\Pi_YUA^{-1}U^*\Pi_Y .
\tag{9.25aq}
\]
The finite adaptive gain is then
\[
 \boxed{
 G_Y:=\mathfrak C_{M,J;X,Y}
 =\langle z,(I+B)^{-1}z\rangle
 =\frac{\tau_{d+1}(Y)}{\tau_d(Y)}.}
\tag{9.25ar}
\]

\begin{theorem}[Exact signed amplitude transfer]
\label{thm:signed-amplitude-transfer}
Let \(\psi\in\ker W^*\) and
\(\beta_\psi=\langle\psi,v\rangle\ne0\).  Define
\[
 g_\psi=(-A^{-1/2}U^*\psi,\psi),
 \qquad
 h_\psi=\frac{(-1)^d\sqrt{\delta_{M,J,X}}}
          {\overline{\beta_\psi}}\,g_\psi .
\tag{9.25as}
\]
If \(\Delta_d\) and \(\Delta_{d+1}\) are the nuisance and augmented
Andreief amplitudes of the column system
\[
 (A^{1/2}e_i,Ue_i),\qquad (0,We_a),\qquad (0,v),
\]
then
\[
 \boxed{
 \int\overline{h_\psi(\omega)}
 \Delta_{d+1}(\omega,\boldsymbol\omega)\,
 \dd\widetilde\nu_Y(\omega)
 =\sqrt{\delta_{M,J,X}}\Delta_d(\boldsymbol\omega).}
\tag{9.25at}
\]
The associated exterior interior-product operator has norm
\[
 \boxed{
 \|T_{h_\psi}\|_{\rm op}^2
 =\delta_{M,J,X}
 \frac{\|\psi\|^2+\|A^{-1/2}U^*\psi\|^2}
  {|\langle\psi,v\rangle|^2}.}
\tag{9.25au}
\]
\end{theorem}

\begin{proof}
The first component of \(g_\psi\) cancels every old-mode column, while
\(W^*\psi=0\) cancels every radical column.  Only the new-mode cofactor
survives the Laplace expansion, giving (9.25at).  Contraction by a
one-particle vector on an exterior power has norm \(\|h_\psi\|\), which
gives (9.25au).
\end{proof}

If \(z=0\), set \(G_Y=Q_k=0\) for all \(k\); no negative pivot can
cross at this cutoff.  Assume henceforth that \(z\ne0\).  The choice
\(\psi=z\) gives the matched-filter certificate
\[
 Q_1=\frac{m_0^2}{m_0+m_1},
 \qquad m_j=\langle z,B^jz\rangle .
\tag{9.25av}
\]
It is sufficient, but it need not detect a crossing already present in
\(G_Y\).

For \(k\ge1\), define
\[
 \mathcal K_k=\operatorname{span}\{z,Bz,\ldots,B^{k-1}z\},
\]
\[
 H_k=[m_{i+j}+m_{i+j+1}]_{i,j=0}^{k-1},\qquad
 b_k=(m_0,\ldots,m_{k-1})^{\mathsf T},\qquad
 Q_k=b_k^*H_k^\dagger b_k .
\tag{9.25aw}
\]

\begin{theorem}[Adaptive Krylov hierarchy]
\label{thm:adaptive-krylov}
One has
\[
 \boxed{
 Q_k=\sup_{0\ne\psi\in\mathcal K_k}
 \frac{|\langle\psi,z\rangle|^2}
  {\langle\psi,(I+B)\psi\rangle},
 \qquad
 Q_1\le Q_2\le\cdots\le G_Y.}
\tag{9.25ax}
\]
If the spectral measure of \(B\) seen from \(z\) has \(r\) support
points, then \(Q_r=G_Y\).
\end{theorem}

\begin{proof}
Writing \(\psi=\sum_{j<k}c_jB^jz\) converts the quotient into
\(|c^*b_k|^2/(c^*H_kc)\).  The finite Riesz formula gives \(Q_k\), and
nestedness gives monotonicity.  If \(c\in\ker H_k\), its Krylov vector is
zero, hence \(c^*b_k=0\); thus \(b_k\in\operatorname{ran}H_k\) and the
pseudoinverse formula is well-defined.  On an \(r\)-point support,
\((1+t)^{-1}\) has a polynomial interpolant of degree at most \(r-1\),
so \((I+B)^{-1}z\in\mathcal K_r\).
\end{proof}

\subsection{Second-level spectral dispersion}
\label{sec:q2-dispersion}

Put
\[
 N_2=m_0m_2-m_1^2,\qquad
 D_2=(m_0+m_1)(m_2+m_3)-(m_1+m_2)^2 .
\tag{9.25ay}
\]
When \(D_2>0\), direct inversion of \(H_2\) gives
\[
 \boxed{
 Q_2=Q_1+\frac{N_2^2}{(m_0+m_1)D_2}.}
\tag{9.25az}
\]
If \(D_2=0\), the spectral measure seen from \(z\) has one support point
and \(Q_2=Q_1=G_Y\).

Let \(d\rho(t)=d\langle z,E_B(t)z\rangle\).  Then
\[
 \boxed{
 N_2=\frac12\iint(s-t)^2\,\dd\rho(s)\dd\rho(t),}
\]
\[
 \boxed{
 D_2=\frac12\iint
 (s-t)^2(1+s)(1+t)\,\dd\rho(s)\dd\rho(t).}
\tag{9.25ba}
\]
Equivalently,
\[
 N_2=\|z\wedge Bz\|^2,\qquad
 D_2=\|(I+B)^{1/2}z\wedge(I+B)^{1/2}Bz\|^2 .
\tag{9.25bb}
\]
Thus, when \(Q_1<\delta_{M,J,X}\),
\[
 \boxed{
 N_2^2>
 (\delta_{M,J,X}-Q_1)(m_0+m_1)D_2}
\tag{9.25bc}
\]
is exactly equivalent to \(Q_2>\delta_{M,J,X}\).

For any two orthonormal observation packets \(\phi_0,\phi_1\), put
\[
 a_i=\langle\phi_i,z\rangle,\qquad
 b_i=\langle\phi_i,Bz\rangle .
\]
The corresponding exterior coordinate gives
\[
 N_2\ge|a_0b_1-a_1b_0|^2.                       \tag{9.25bd}
\]
Since \(D_2\le(1+\|B\|)^2N_2\), the signed two-packet condition
\[
 \boxed{
 |a_0b_1-a_1b_0|^2>
 (\delta_{M,J,X}-Q_1)(m_0+m_1)(1+\|B\|)^2}
\tag{9.25be}
\]
implies the finite crossing.

The literal ordinary-prime form uses one common radical and old-mode
regression.  Put
\[
 \zeta^*=-(W^*W)^{-1}W^*v,\qquad
 Z=-(W^*W)^{-1} W^*U,
\]
\[
 \widetilde q=q_{M,J,X}^*+\Psi\zeta^*,\qquad
 \widetilde\Phi=\Phi+\Psi Z,
\]
and, in the individual prime-power channels,
\[
 x_n=\mathcal D_n\widetilde q,\qquad
 F_n=\mathcal D_n\widetilde\Phi A^{-1/2},\qquad
 \xi=\sum_{\substack{n=p^k\\X<n\le Y}}F_n^*x_n .
\tag{9.25bf}
\]
Then
\[
 \boxed{
 N_2=\frac12
 \sum_{\substack{n=p^k,\ \ell=q^a\\X<n,\ell\le Y}}\iint
 \left|x_n(s)(F_\ell\xi)(t)
 -(F_n\xi)(s)x_\ell(t)\right|^2\,\dd s\,\dd t.}
\tag{9.25bg}
\]
All von Mangoldt weights and theta phases therefore remain coupled until
the final modulus square.

\begin{corollary}[Exact hierarchy formulation of the frontier]
\label{cor:krylov-frontier}
For every fixed negative anti-shorted row,
\[
 \boxed{
 \mathfrak g_{M,J,X}>\delta_{M,J,X}
 \Longleftrightarrow
 \exists\,Y<\infty,\ \exists\,k<\infty:
 Q_k(Y)>\delta_{M,J,X}.}
\tag{9.25bh}
\]
\end{corollary}

\begin{proof}
The forward implication follows from augmented determinant exhaustion,
followed by finite Krylov exactness.  Conversely,
\(Q_k(Y)\le G_Y\le\mathfrak g_{M,J,X}\).
\end{proof}

\begin{remark}[Level two is a sufficient gate]
The depth in Corollary~\ref{cor:krylov-frontier} may depend on the finite
row.  The dispersion inequality (9.25bc) is the first explicit
strengthening of MF, not a replacement for the exact surplus.  For
example,
\[
 B=\operatorname{diag}(1,2,3),\qquad
 z=\sqrt{12/13}(1,1,1),\qquad
 \delta=649/650
\]
gives \(G_Y=1>\delta\), but \(Q_2=648/650<\delta\).  Failure at level
two is therefore inconclusive.  If \(G_Y>\delta\), finite Krylov
exactness guarantees that some higher row-dependent level detects the
crossing.
\end{remark}

\begin{problem}[Physical strict surplus]
\label{prob:physical-strict-surplus}
For every negative anti-shorted row required in the cofinal exhaustion,
prove, with the literal theta translations and ordinary von Mangoldt
weights,
\[
 \boxed{\mathfrak g_{M,J,X}>\delta_{M,J,X}.}
\tag{9.25bi}
\]
By Theorem~\ref{thm:augmented-exhaustion}, this is equivalent to one
finite inequality (9.25ao), and by
Corollary~\ref{cor:krylov-frontier} to one finite Krylov certificate of
row-dependent depth.  No localization, atomwise minimization, or
infinite-selection loss remains in this formulation.
\end{problem}

The mean-periodic constraint must be applied before the finite radical
correction.  Indeed, if
\(q\in(\mathbf1\oplus\mathcal R)^\perp\) and
\(r\in\mathfrak R_J\), then
\[
 (h(q+r))*K=0
 \quad\Longleftrightarrow\quad r=0,
\]
because the mean-periodic complement intersects the complete radical only
at zero.  Thus \(q_{M,J,X}^*+d\) is mean-periodic for
\(d\in V_{M-1}\), whereas a joint saddle vector with a nonzero radical
correction is not.  The correction must remain in the determinant, but the
convolution equation may not be imposed on it.

\subsection{Heat-core exhaustion of the radical complement}
\label{subsec:heat-core}

The elementary zero modes provide the explicit orbit coordinate used in
the finite calculations above.  Their completeness in the weighted form
topology is a separate spectral-synthesis statement.  It is not needed for
the global passage.  We now construct a form core directly from the closed
ordinary-prime--Gamma form.

Let
\[
 \mathscr C=(\mathbf1\oplus\mathcal R)^\perp,
 \qquad
 \mathcal Q_{\mathscr C}=D(\mathscr E_K)\cap\mathscr C.
\tag{9.25bj}
\]
The complete form is closed.  Indeed, on its maximal domain it is the
squared norm of the difference operator
\[
\begin{aligned}
 \mathcal Dq={}&
 \left(
 \sqrt{g(u)K(x)K(x-u)}\{q(x)-q(x-u)\}
 \right)_{(u,x)\in(0,\infty)\times\mathbb R}
 \\
 &\oplus
 \left(
 \sqrt{\frac{\Lambda(n)}{\sqrt n}K(x)K(x-\log n)}
 \{q(x)-q(x-\log n)\}
 \right)_{n=p^a},
\end{aligned}
\tag{9.25bk}
\]
from \(L^2_{\rm even}(\mu_K)\) to the corresponding Hilbert direct sum.
If \(q_\ell\to q\) and \(\mathcal Dq_\ell\to Y\), one subsequence converges
almost everywhere with respect to Lebesgue measure because \(d\mu_K\) has
a strictly positive density.  Fubini gives simultaneous convergence at
\(x-u\), and an almost-everywhere subsequence in the disjoint-union output
space identifies \(Y=\mathcal Dq\).  Thus \(\mathcal D\) is closed.
Moreover \(C_{c,\rm even}^\infty(\mathbb R)\subset D(\mathcal D)\): near
\(u=0\), \(g(u)=(2u)^{-1}+O(1)\) and the squared difference is \(O(u^2)\),
while for large \(u\) the theta overlap is \(O_q(e^{-c_qe^{2u}})\).
At \(u=\log n\) this is summable against
\(\Lambda(n)/\sqrt n\).  Hence the form is densely defined and closed.
Let \(L\ge0\) be its associated self-adjoint operator.

Equation (9.15), interpreted by the representation theorem, gives the
strong domain identities
\[
 L\mathbf1=0,
 \qquad
 r_j\in D(L),\qquad Lr_j=\frac12r_j.
\tag{9.25bl}
\]
Therefore \(\mathbf1\oplus\mathcal R\) and \(\mathscr C\) reduce \(L\).
Put
\[
 S_{\mathscr C}=L|_{\mathscr C}+\frac12I\ge\frac12I.
\tag{9.25bm}
\]

Choose a sequence \((g_j)_{j\ge1}\subset\mathscr C\) whose algebraic span
is Hilbert-dense and define
\[
 V_M^{\rm heat}
 =\operatorname{span}\{e^{-S_{\mathscr C}/k}g_j:1\le j,k\le M\}.
\tag{9.25bn}
\]

\begin{theorem}[Heat-core Galerkin exhaustion]
\label{thm:heat-core}
The spaces \(V_M^{\rm heat}\) are nested finite-dimensional subspaces of
\(\mathcal Q_{\mathscr C}\), every one of their vectors belongs to
\(\cap_{m\ge0}D(S_{\mathscr C}^m)\), and
\[
 \boxed{
 \overline{\bigcup_MV_M^{\rm heat}}^{\,
 (\frac12\|\cdot\|_{\mu_K}^2+\mathscr E_K)^{1/2}}
 =\mathcal Q_{\mathscr C}.}
\tag{9.25bo}
\]
For every fixed finite-dimensional
\(E\subset D(\mathscr E_K)\), the omitted Euler tail satisfies
\[
 \boxed{
 \left\|N_E^{-1/2}(\mathcal T_X)_EN_E^{-1/2}\right\|
 \longrightarrow0.}
\tag{9.25bp}
\]
Finally, if a finite omitted block contains two distinct primes \(p,r\),
its Gram form is positive definite on every finite centered space
\(E\oplus\mathfrak R_J\) with \(E\subset\mathscr C\).
\end{theorem}

\begin{proof}
Spectral calculus gives, for \(t>0\),
\[
 \|S_{\mathscr C}^{m}e^{-tS_{\mathscr C}}\|
 \le\left(\frac{m}{et}\right)^m,
 \qquad
 \|S_{\mathscr C}^{1/2}e^{-tS_{\mathscr C}}\|
 \le(2et)^{-1/2}.
\tag{9.25bq}
\]
For \(q\in D(S_{\mathscr C}^{1/2})\), dominated convergence in its spectral
measure gives
\[
 \|S_{\mathscr C}^{1/2}(e^{-tS_{\mathscr C}}-I)q\|\longrightarrow0.
\tag{9.25br}
\]
Given \(\varepsilon>0\), first choose \(k\) so that (9.25br), with
\(t=1/k\), is smaller than \(\varepsilon/2\).  Approximate \(q\) in the
Hilbert norm by a finite linear combination \(g\) of the \(g_j\), with
\(\|q-g\|<(\varepsilon/2)\sqrt{2et}\).  Equations (9.25bq)--(9.25br)
then give
\[
 \|S_{\mathscr C}^{1/2}(q-e^{-tS_{\mathscr C}}g)\|<\varepsilon,
\]
which proves (9.25bo).

For (9.25bp), choose a basis \(e_1,\ldots,e_d\) of \(E\).  Positivity and
monotone convergence give
\(\mathcal T_X(e_i,e_i)\downarrow0\), while Cauchy--Schwarz for the tail
form gives
\[
 |\mathcal T_X(e_i,e_j)|^2
 \le\mathcal T_X(e_i,e_i)\mathcal T_X(e_j,e_j).
\]
Thus every entry, and hence the normalized finite matrix, tends to zero.

For the last assertion, vanishing of the two prime atoms would give the
two periods \(\log p\) and \(\log r\).  Their quotient is irrational,
since a rational relation would imply \(p^a=r^b\).  Strong continuity of
translations in \(L^2_{\rm loc}\) then makes every real number a period,
so the vector is constant almost everywhere.  Centeredness makes it zero.
The finite Gram matrix is therefore positive definite.
\end{proof}

The hybrid spaces
\[
 V_M^{\rm hyb}=V_M^{\rm heat}+V_M
\tag{9.25bs}
\]
are also a form core.  They retain every explicit zero-orbit diagnostic
used above, without asserting that the canonical zero modes alone are
dense.  One-channel observability and the local scalar Christoffel bounds
remain statements about the canonical summand.  In the global closure,
finite tail blocks must contain two distinct primes, and the physical
surplus must be proved on the heat or hybrid rows.

\subsection{The compensated frontier and closure}

The strict frontier (9.12) excludes a completed threshold null vector and
can therefore be stronger than mere nonnegativity.  The exact
vanishing-tolerance formulation is the following.  Choose
\[
 M\longmapsto(J(M),X(M),\varepsilon_M),
\quad J(M),X(M)\to\infty,\quad\varepsilon_M\downarrow0,
\tag{9.26}
\]
and require
\[
 \boxed{
 n_-\!\left[
 \mathbf A_{M,J(M),X(M)}
 +\varepsilon_M
 \begin{pmatrix}\mathbf N_M&0\\0&0\end{pmatrix}
 \right]=J(M).}
\tag{9.27}
\]
Equivalently,
\[
 \boxed{
 \mathbf T^{VV}
 -\mathbf T^{VR}(\mathbf T^{RR})^{-1}\mathbf T^{RV}
 \preceq\mathbf H_M+\varepsilon_M\mathbf N_M.}
\tag{9.28}
\]
On the canonical zero-mode summand the cutoff part of (9.28) has the
explicit bound \(C_Me^{-cX(M)}\mathbf N_M\).  On every heat or hybrid row,
Theorem~\ref{thm:heat-core} gives the qualitative bound needed for the
diagonal schedule: after the finite row is fixed, \(X(M)\) may be enlarged
until the normalized tail is at most \(\varepsilon_M\).  The remaining
content is the joint signed prime--Gamma floor.

\begin{theorem}[Cofinal inertia closure]
\label{thm:cofinal-inertia-closure}
Let \(V_M\) be either the heat spaces (9.25bn) or the hybrid spaces
(9.25bs), and assume that (9.27) holds.
Then
\[
 \mathcal A_\infty(q,q)\ge0
\tag{9.29}
\]
on that complement.  Under the moving/model identification this gives
the source-side Mosco lower bound, and hence RH.
\end{theorem}

\begin{proof}
First synchronize the inertia and tail cutoffs.  If (9.27) holds at
\(\widehat X(M)\), replace it by any
\[
 X(M)\ge\max\{\widehat X(M),X_{\rm tail}(M,J(M),\varepsilon_M)\}.
\]
Increasing the cutoff adds a positive Gram matrix, so the negative index
cannot increase.  It cannot fall below \(J(M)\), because the shifted form
on every nonzero \(r\in\mathfrak R_{J(M)}\) is
\[
 \mathcal A_{X(M)}(r,r)=-\mathcal T_{X(M)}(r,r)<0;
\]
the remaining tail contains two distinct primes and
Theorem~\ref{thm:heat-core} applies.  Thus (9.27) persists at the
synchronized cutoff.

Now fix \(q\in V_m\).  For every \(M\ge m\), (9.27), the Schur-complement
formula, and the inclusion of \(q\) in the larger mode block give
\[
 \mathcal A_\infty(q,q)
 \ge-\varepsilon_M\|q\|_{\mu_K}^2.
\]
Let \(M\to\infty\).  Then \(\varepsilon_M\to0\), so
\(\mathcal A_\infty(q,q)\ge0\) on \(\cup_MV_M\).
Theorem~\ref{thm:heat-core} and continuity in the closed form norm extend
the inequality to the complete centered radical complement.  The radical
has value zero by (9.15), while constants have zero value in
\(\mathscr E_K-\frac12\operatorname{Var}_{\mu_K}\).  Hence the complete
centered, equivalently constant-quotiented, Weil form is nonnegative.
Theorem~\ref{thm:weil} gives RH.  Equivalently, in the CCM
branch coordinate, Theorems~\ref{thm:even-rayleigh-residual} and
\ref{thm:cross-residual-vanishing} remove the already controlled residual
pieces and (9.27) supplies the source liminf in
Problem~\ref{prob:source-mosco}.
\end{proof}

\begin{center}
\begin{tabular}{>{\raggedright\arraybackslash}p{0.34\textwidth}
                >{\raggedright\arraybackslash}p{0.47\textwidth}}
\toprule
Component & Status in this proof chain\\
\midrule
CCM finite self-adjoint quotient & Published finite package.\\
Connes prolate critical-line model & Published finite/model package.\\
Suzuki localized operator & Published localized realization.\\
Even Rayleigh residual & Proved:
\(\lvert\mathscr R_L^+\rvert\ll\lambda^7d_4\).\\
Direct cross residual & Proved:
\(\|b_L^+\|\ll\lambda^4\sqrt{(1+L)d_4}\to0\).\\
Literal Euler-head approximation & Explicit canonical rate proved by
(9.7)--(9.8); qualitative normalized convergence on every heat/hybrid row
proved by Theorem~\ref{thm:heat-core}.\\
Finite radical inertia algebra & Proved by (9.16)--(9.22).\\
Large-prime channel geometry & Proved by (9.23)--(9.25).\\
Finite one-channel observability & Proved on the canonical zero-mode rows by
Theorem~\ref{thm:one-channel-observability}; finite rank and aliasing are
removed there.  General heat rows use two-prime blocks.\\
Completed scalar innovation & Exact recursion proved by
Theorem~\ref{thm:scalar-innovation}.\\
Prime-by-prime staircase update & Exact positive Kalman--Schur increment
proved in (9.25j)--(9.25k), with the literal \(\Lambda(p^k)\) weights.\\
Local and cumulative gain & Midpoint Christoffel detection and the exact
block telescope are proved in Proposition~\ref{prop:local-christoffel}
and Proposition~\ref{prop:cumulative-gain}.\\
Radical-conditioned gain & The exact projected residual, resolvent gain,
and matched-filter criterion are proved in
Theorem~\ref{thm:radical-conditioned-gain}.\\
Augmented finite determinant & Strict positivity, determinant formula,
monotone exhaustion, and finite selection are proved in
Theorem~\ref{thm:augmented-exhaustion}.\\
Signed amplitude transfer & The exact exterior interior-product identity
and its norm
are proved in Theorem~\ref{thm:signed-amplitude-transfer}.\\
Adaptive Krylov recovery & Monotonicity and finite-row exactness are
proved in Theorem~\ref{thm:adaptive-krylov}.\\
Second-level dispersion gate & Formula (9.25az), its spectral wedge
identity, and the sufficient test (9.25bc) are proved; the test is not
asserted uniformly for every physical row.\\
Physical strict surplus (9.25bi) & Remaining force-bearing arithmetic
inequality; it must hold on the heat or hybrid form-core rows, with no
localization or infinite-selection loss.\\
Joint compensated frontier (9.27) & Follows row by row from the strict
surplus; its vanishing-tolerance form permits threshold equality.\\
Prime--Gamma form-core exhaustion & Proved by
Theorem~\ref{thm:heat-core}; canonical zero-mode-only synthesis is not
asserted or required.\\
\bottomrule
\end{tabular}
\end{center}

\begin{remark}[Strict and compensated targets]
The strict statement requested by the finite-energy-level picture is
\[
 \boxed{X_*(M,J(M))<\infty\quad\text{for every }M.}
\tag{9.30}
\]
By Theorem~\ref{thm:exact-finite-repair}, its force-bearing content is a
strict completed finite-mode gap.  The compensated statement (9.27) is
the threshold-stable version needed for RH.  A proof of (9.30) on a
form-core exhaustion proves (9.27) immediately; a proof of (9.27) need
not exclude additional threshold null states at finite \(M\).
\end{remark}

\section{Tate boundary geometry and the global polarization audit}
\label{sec:polarization-audit}

The finite-inertia analysis above isolates the missing sign but does not
explain where that sign could come from.  This section records a second,
source-level construction.  It starts from finite root graphs at the
ordinary primes, constructs the Gamma and polar determinant at infinity,
and joins them before completion.  The construction produces the correct
local masses, parity, boundary cancellation and scale covariance.  It also
permits a decisive audit of the ordinary Hilbert completion.

The distinction maintained throughout is
\[
 \boxed{
 \text{positive local/source polarization}
 \ \not\Longrightarrow\ 
 \text{positive faithful polarization of }H^1_{\rm CCM}.}
 \tag{10.1}
\]
The left side will be constructed below.  The most natural absolutely
continuous realization of the arrow in (10.1) will be disproved.

\subsection{Finite root geometry and the literal von Mangoldt mass}

For \(N\ge2\), let \(R_N=\mathbb Z/N\mathbb Z\), let \(T_N\) be
translation by one on \(\ell^2(R_N)\), and put
\[
 \Delta_N=2I-T_N-T_N^*.
 \tag{10.2}
\]
Its form is the positive graph energy
\[
 \langle f,\Delta_Nf\rangle
 =\sum_{j\in R_N}|f(j+1)-f(j)|^2.
 \tag{10.3}
\]
The kernel consists of the constants.  We write \(\det'\Delta_N\) for
the product of the nonzero eigenvalues.

\begin{theorem}[Root torsion--incidence identity]
\label{thm:root-torsion}
For every \(N\ge2\),
\[
 \det'\Delta_N=N^2.
 \tag{10.4}
\]
Fix \(m\ge2\), a prime \(p\), and \(k\ge1\).  Then
\[
 \frac12\log
 \frac{\det'\Delta_{mp^k}}{\det'\Delta_{mp^{k-1}}}
 =\log p.
 \tag{10.5}
\]
If \(H_{m,k}=p^kR_m\subset R_{mp^k}\), and
\[
 \omega_{mp^k}=(mp^k)^{-1/2}\mathbf1_{R_{mp^k}},
 \qquad
 \eta_{m,k}=m^{-1/2}\mathbf1_{H_{m,k}},
\]
then
\[
 \langle\omega_{mp^k},\eta_{m,k}\rangle=p^{-k/2}.
 \tag{10.6}
\]
Consequently
\[
 \boxed{
 \left(\frac12\log
 \frac{\det'\Delta_{mp^k}}{\det'\Delta_{mp^{k-1}}}\right)
 \langle\omega_{mp^k},\eta_{m,k}\rangle
 =\frac{\Lambda(p^k)}{\sqrt{p^k}}.}
 \tag{10.7}
\]
\end{theorem}

\begin{proof}
The Fourier characters of \(R_N\) diagonalize \(\Delta_N\), with
eigenvalues \(|1-e^{2\pi ij/N}|^2\).  Hence
\[
 \det'\Delta_N
 =\left|\prod_{j=1}^{N-1}(1-e^{2\pi ij/N})\right|^2=N^2,
\]
which gives (10.4) and (10.5).  The subgroup \(H_{m,k}\) has \(m\)
elements, so the normalized overlap is
\(m/(\sqrt{mp^k}\sqrt m)=p^{-k/2}\).  Multiplying by (10.5)
proves (10.7).
\end{proof}

The two factors in (10.7) have different meanings.  The logarithm is a
relative determinant increment along the root tower, whereas
\(p^{-k/2}\) is a normalized incidence number.  The distinction prevents
the incorrect replacement of the weight of the layer \(p^k\) by
\(k\log p\).

The finite root levels are polarized before any global sign is considered.
On \(H_N^0=\mathbf1^\perp\subset\ell^2(R_N)\), set
\[
 \begin{aligned}
 J_N(x,y)&=(-y,x),\\
 g_N((x,y),(u,v))
 &=\langle\Delta_Nx,u\rangle+\langle\Delta_Ny,v\rangle,\\
 \Omega_N(z,w)&=-g_N(z,J_Nw).
 \end{aligned}
 \tag{10.8}
\]
Then \(g_N\) is positive definite, \(J_N^2=-I\), and \(\Omega_N\) is
alternating and nondegenerate.  If \(M\mid N\), normalized pullback along
\(R_N\to R_M\) is a polarized isometry and the determinant on the new-root
complement is \((N/M)^2\).  Thus (10.8) is functorial under the root
coverings which produce (10.5).

\subsection{Poisson holonomy, heat lift and local unitary triviality}

The coefficient \(p^{-k/2}\) also has an exact positive realization as a
unitary moment.  For \(0<r<1\), let
\[
 \dd\mu_r(\theta)
 =\frac{1-r^2}{1-2r\cos\theta+r^2}\frac{\dd\theta}{2\pi}.
\tag{10.8a}
\]
The Fourier expansion of the Poisson kernel gives
\[
 \int_{-\pi}^{\pi}e^{ik\theta}\dd\mu_r(\theta)=r^{|k|}.
\tag{10.8b}
\]
Taking \(r=p^{-1/2}\) realizes every winding of the prime tower without
analytic continuation.

Let \(A_{p,\theta}\) be the nonnegative Laplacian on the circle of length
\(\ell_p=\log p\), with boundary condition
\(f(x+\ell_p)=e^{i\theta}f(x)\).  The method-of-images kernel is
\[
 K_{p,\theta,t}(x,y)
 =\frac1{\sqrt{4\pi t}}
  \sum_{k\in\mathbb Z}
  e^{-(x-y+k\ell_p)^2/(4t)}e^{ik\theta}.
\tag{10.8c}
\]
It satisfies the exact heat composition law before any winding sector is
removed.  Integrating its trace against (10.9) yields
\[
 \boxed{
 \tau_p(e^{-tA_p})
 =\frac{\ell_p}{\sqrt{4\pi t}}
  \left(1+2\sum_{k\ge1}p^{-k/2}
  e^{-(k\ell_p)^2/(4t)}\right).}
\tag{10.8d}
\]
The nonzero-winding relative trace is therefore exactly
\[
 \tau_p(e^{-tA_p})-\frac{\log p}{\sqrt{4\pi t}}
 =\frac{\log p}{\sqrt{\pi t}}
  \sum_{k\ge1}p^{-k/2}e^{-(k\log p)^2/(4t)}.
\tag{10.8e}
\]
This is the complete Gaussian prime tower in the heat explicit formula.

The local realization has a rigidity which is important globally.  The
Poisson density is the modulus square of the outer function
\[
 a_r(e^{i\theta})=\frac{\sqrt{1-r^2}}{1-re^{i\theta}}.
\tag{10.8f}
\]
Multiplication by \(a_r\) is a unitary map
\(L^2(\mathbb T,\mu_r)\to L^2(\mathbb T,\dd\theta/2\pi)\) and commutes
with every decomposable fiber operator.  Hence the weighted and unweighted
local Laplacians are unitarily equivalent.  The prime changes the normal
weight used to read the heat diagonal; it does not create a new local
spectrum.  Moreover the winding-zero sector cannot be deleted while
preserving composition, because two nonzero windings can add to zero.
Thus the finite quantity in (10.13) is a relative semifinite trace, not the
ordinary trace of a reduced semigroup.

There is an equivalent precision-space formulation.  On
\(\ell^2(\mathbb N_0)\), let
\[
 (K_px)_{h}=\sum_{k\ge0}p^{-|h-k|/2}x_k.
\]
With \(a=p^{-1/2}\), define
\[
 (B_ax)_0=x_0,
 \qquad
 (B_ax)_{h+1}=\frac{x_{h+1}-ax_h}{\sqrt{1-a^2}}.
\tag{10.8g}
\]
Direct triangular inversion gives
\[
 \boxed{
 K_p^{-1}=B_a^*B_a,
 \qquad
 \langle K_p^{-1}x,x\rangle
 =|x_0|^2+\frac1{1-p^{-1}}
  \sum_{h\ge0}|x_{h+1}-p^{-1/2}x_h|^2.}
\tag{10.8h}
\]
Thus the Poisson coefficient is also the Green kernel of a positive local
first-order arithmetic connection.  The local positivity is exact, while
\(\sum_p p^{-1}=\infty\) prevents the infinite product of these changes of
coordinates from being a bounded perturbation of a reference product.
The matched Gamma--Euler boundary must therefore be formed before the
global norm or trace is taken.

\subsection{The Gamma spin and the polar determinant}

On \(\mathcal H_\Gamma=\ell^2(\mathbb N_0)\), define the positive operator
\[
 N_\Gamma e_m=\left(2m+\frac12\right)e_m.
\tag{10.9}
\]
Its heat trace is
\[
 \Tr e^{-uN_\Gamma}=\frac{e^{-u/2}}{1-e^{-2u}}.
\tag{10.10}
\]
Zeta regularization of \(N_\Gamma+s-1/2\) gives the archimedean
determinant directly from this source operator.

\begin{theorem}[Gamma--polar determinant]
\label{thm:gamma-polar-det}
For \(\Re s>0\), with meromorphic continuation thereafter,
\[
 \det_\zeta(N_\Gamma+s-1/2)
 =\frac{\sqrt{2\pi}\,2^{1/2-s/2}}{\Gamma(s/2)}.
\tag{10.11}
\]
Let \(N_{\rm triv}\) have eigenvalues \(-1/2\) and \(1/2\).  Then
\[
 \boxed{
 \frac12s(s-1)\pi^{-s/2}\Gamma(s/2)
 =\sqrt\pi(2\pi)^{-s/2}
 \frac{\det((s-1/2)I-N_{\rm triv})}
      {\det_\zeta(N_\Gamma+s-1/2)}.}
\tag{10.12}
\]
\end{theorem}

\begin{proof}
The spectral zeta function is
\[
 \Tr(N_\Gamma+s-1/2)^{-w}
 =2^{-w}\zeta_H(w,s/2).
\]
Use
\(\zeta_H(0,a)=1/2-a\) and
\(\zeta_H'(0,a)=\log\Gamma(a)-\frac12\log(2\pi)\) to obtain
(10.11).  Since
\(\det((s-1/2)I-N_{\rm triv})=s(s-1)\), substitution gives
(10.12).
\end{proof}

The positive degree-one Gamma page is the real double of the form domain
of \(N_\Gamma^{1/2}\), with metric and complex structure
\[
 \begin{aligned}
 g_\Gamma((x,y),(u,v))
 &=\langle N_\Gamma^{1/2}x,N_\Gamma^{1/2}u\rangle
  +\langle N_\Gamma^{1/2}y,N_\Gamma^{1/2}v\rangle,\\
 J_\Gamma(x,y)&=(-y,x).
 \end{aligned}
\tag{10.13}
\]
The polar plane remains in degrees zero and two.  It is not inserted as a
positive degree-one summand.  This preserves the Lefschetz parity
\(H^0-H^1+H^2\), which is responsible for the sign of the prime terms in
the explicit formula and is central to the geometric interpretation in
\S\ref{subsec:open-programs}.

\subsection{Matched boundary cancellation and the full finite part}

Let \(\mathscr K\) denote the common Cauchy coefficient space.  For each
prime power set
\[
 w_{p,k}=\frac{\log p}{p^k}.
 \tag{10.14}
\]
The restricted-product boundary is decomposed into a common component
\(F\in\mathscr K\) and local deviations \(r=(r_{p,k})\), where
\[
 \sum_{p,k}w_{p,k}\|r_{p,k}\|
 +\left(\sum_{p,k}w_{p,k}\|r_{p,k}\|^2\right)^{1/2}<\infty.
 \tag{10.15}
\]
Thus the local row is \(F_{p,k}=F+r_{p,k}\).  For \(s>1/2\), put
\[
 C_s=\sum_p(\log p)p^{-2s},
 \qquad
 \kappa_\infty=\gamma+\sum_p\sum_{k\ge2}w_{p,k}.
 \tag{10.16}
\]
The ordinary prime number theorem, in its logarithmic-derivative form at
the pole, gives the normalized finite-part identity
\[
 \operatorname{FP}_{s\downarrow1/2}C_s=-\kappa_\infty.
 \tag{10.17}
\]

For \(x=(F,r)\), \(y=(G,z)\), consider the jointly regularized form
\[
 \begin{aligned}
 \mathfrak Q_s(x,y)
 &=\sum_p(\log p)p^{-2s}
   \langle F+r_{p,1},G+z_{p,1}\rangle\\
 &\quad+\gamma\langle F,G\rangle
 +\sum_p\sum_{k\ge2}w_{p,k}
   \langle F+r_{p,k},G+z_{p,k}\rangle.
 \end{aligned}
 \tag{10.18}
\]

\begin{theorem}[Exact matched finite part]
\label{thm:matched-finite-part}
The finite part of (10.18) exists and equals
\[
 \boxed{
 \operatorname{FP}_{s\downarrow1/2}\mathfrak Q_s(x,y)
 =\sum_{p,k}w_{p,k}
 \bigl(
 \langle r_{p,k},G\rangle
 +\langle F,z_{p,k}\rangle
 +\langle r_{p,k},z_{p,k}\rangle
 \bigr).}
 \tag{10.19}
\]
The series is absolutely convergent and defines a continuous form on the
space in (10.15).  In particular the common generic diagonal has zero
finite part.
\end{theorem}

\begin{proof}
After expanding (10.18), the coefficient of \(\langle F,G\rangle\) is
\(C_s+\kappa_\infty\), whose finite part is zero by (10.17).  Every
residual coefficient converges to \(w_{p,k}\).  The first summand in
(10.15) controls the common--residual terms and weighted
Cauchy--Schwarz controls the residual--residual term, so dominated
convergence gives (10.19).
\end{proof}

Theorem~\ref{thm:matched-finite-part} is stronger than cancellation of
two scalar divergences.  It retains every independent local deviation with
its literal von Mangoldt weight.  Hence Gamma, the pole and the prime
layers are joined before taking an absolute square or a quotient.

\subsection{Covariant mixing of all prime-power charges}

The local row of charge \(\ell_{p,k}=k\log p\) has normalized scale
generator \(A-\ell_{p,k}\).  Adding rows with different charges without
transporting them is not equivariant.  The required transport is explicit.
In the spectral realization
\[
 \mathscr K_C=L^2(\mathbb R,w_C(\gamma)\dd\gamma),
 \qquad
 w_C(\gamma)=\frac1{2\pi(\gamma^2+1/4)},
 \qquad AF(\gamma)=\gamma F(\gamma),
 \tag{10.20}
\]
define
\[
 (S_\ell F)(\gamma)
 =\left(\frac{w_C(\gamma+\ell)}{w_C(\gamma)}\right)^{1/2}
  F(\gamma+\ell).
 \tag{10.21}
\]
Then \((S_\ell)_{\ell\in\mathbb R}\) is a strongly continuous unitary
group and
\[
 S_\ell(A-\ell I)=AS_\ell,
 \qquad
 S_\ell^*f(A)S_\ell=f(A-\ell I).
 \tag{10.22}
\]

For \(I_X=\{(p,k):p^k\le X\}\), put
\[
 c_p=\frac{2\pi}{\log p},
 \qquad
 \alpha_{p,k}=\sqrt{\frac{\log p}{c_p}}p^{-k/2}.
 \tag{10.23}
\]
The covariant Tate boundary is
\[
 \begin{aligned}
 R_X^{\rm cov}v
 &=\sum_{p,k}\alpha_{p,k}S_{\ell_{p,k}}y_{p,k},\\
 R_X^{\rm cov}J_Xv
 &=\sum_{p,k}c_p\alpha_{p,k}S_{\ell_{p,k}}x_{p,k}.
 \end{aligned}
 \tag{10.24}
\]
It intertwines every charged source action with the one action generated
by \(A\).  Its adjoint generic plane has Gram mass
\[
 C_XI,
 \qquad
 C_X=\sum_{p^k\le X}\frac{\log p}{p^k}.
 \tag{10.25}
\]

Let the complete Gamma compliance be
\[
 K_\Gamma=\kappa_\infty I+m_\Gamma(A)succeq\kappa_\infty I.
 \tag{10.26}
\]
Shorting only the common archimedean boundary gives the positive metric
\[
 \begin{aligned}
 g_X^{\rm cov}(v,w)
 &=g_X(v,w)\\
 &\quad+\langle K_\Gamma^{-1}R_X^{\rm cov}J_Xv,
                         R_X^{\rm cov}J_Xw\rangle\\
 &\quad+\langle K_\Gamma^{-1}R_X^{\rm cov}v,
                         R_X^{\rm cov}w\rangle.
 \end{aligned}
 \tag{10.27}
\]
The diagonal row is
\[
 S_{\ell_{p,k}}^*K_\Gamma^{-1}S_{\ell_{p,k}}
 =\bigl(\kappa_\infty+m_\Gamma(A-\ell_{p,k})\bigr)^{-1},
 \tag{10.28}
\]
while two different rows retain the cross compliance
\[
 S_{\ell_i}^*K_\Gamma^{-1}S_{\ell_j}.
 \tag{10.29}
\]
Thus all prime rows use one common regression and all off-diagonal phases
survive.  Equations (10.19) and (10.22) show that the matched finite-part
identity is unchanged in the decharged frame.

\subsection{The global polarization target}

Let \(H^1_{\rm CCM}\) denote the separated resonant degree one supplied by
the CCM cyclic construction, with weight-one scale action
\(\vartheta_t\) and descended alternating form \(\Omega\).  The missing
global object can now be stated independently of the finite matrix model.

\begin{definition}[Global CCM polarization]
\label{def:global-polarization}
A global polarization is a real-linear operator \(J\) on a dense
scale-invariant core of \(H^1_{\rm CCM}\) such that
\[
 J^2=-I,
 \qquad
 J\vartheta_t=\vartheta_tJ,
 \qquad
 g(u,v):=\Omega(u,Jv)
 \tag{10.30}
\]
extends to a faithful positive Hilbert metric.  The normalized action
\(U_t=e^{-t/2}\vartheta_t\) must extend unitarily to its completion.
\end{definition}

This is the infinite-dimensional analogue of the polarized degree-one
page in the proof of the Riemann hypothesis for curves over finite fields,
but the required positivity there is supplied by correspondences and the
Hodge index theorem on a square, rather than by a formal Hilbert
completion.  Deninger's dynamical-cohomology program gives the same
weight-one template~\cite{Deninger1998}.

\begin{proposition}[Exact force of the polarization]
\label{prop:polarization-rh}
A global CCM polarization in the sense of
Definition~\ref{def:global-polarization} implies RH.
\end{proposition}

\begin{proof}
If a resonant class \(u_\rho\ne0\) has
\(\vartheta_tu_\rho=e^{t\rho}u_\rho\), then
\[
 U_tu_\rho=e^{t(\rho-1/2)}u_\rho.
\]
Unitarity for every real \(t\) forces
\(e^{t(\Re\rho-1/2)}=1\), hence \(\Re\rho=1/2\).  The CCM resonant
divisor is the nontrivial zero divisor of the completed zeta function, so
all its points lie on the critical line.
\end{proof}

Proposition~\ref{prop:polarization-rh} explains why positivity of the local
pages cannot be promoted by a formal completion: the promotion itself
contains the full missing sign.

\subsection{The absolutely continuous completion no-go}

The covariant system (10.20)--(10.29) has an ordinary Hilbert direct-limit
completion.  It is natural, positive and source-defined, but it has the
wrong spectral type.  The following operator-valued statement gives the
precise obstruction.

\begin{lemma}[Operator-valued dense-range criterion]
\label{lem:operator-dense-range}
Let \((X,\mu)\) be sigma finite, let \(\mathfrak h\) be separable, and let
\(M(x)\) be a measurable field of closed densely defined operators.  If
\(\mathcal M\) is the associated maximal decomposable operator on
\(L^2(X,\mu;\mathfrak h)\), then
\[
 \boxed{
 \overline{\operatorname{Ran}\mathcal M}
 =L^2(X,\mu;\mathfrak h)
 \quad\Longleftrightarrow\quad
 \ker M(x)^*=\{0\}\ \text{for almost every }x.}
 \tag{10.31}
\]
If \(M(x)^*M(x)\succeq\kappa I\) almost everywhere for one
\(\kappa>0\), then \(\mathcal M\) has closed range and is bounded below.
\end{lemma}

\begin{proof}
The adjoint of the maximal decomposable operator has fibers \(M(x)^*\).
Therefore
\[
 (\overline{\operatorname{Ran}\mathcal M})^\perp
 =\ker\mathcal M^*
 =\int_X^\oplus\ker M(x)^*\dd\mu(x),
\]
which proves (10.31).  The lower-bound assertion is the closed-range
theorem applied fiberwise and then integrated.
\end{proof}

Apply the lemma to the charged Gamma operator
\[
 K_{\Gamma,Q}
 =\kappa_\infty I+\bigoplus_{q\in Q}m_\Gamma(A-\log q).
 \tag{10.32}
\]
It is uniformly bounded below by \(\kappa_\infty I\).  If the scalar term
is omitted, the fiber kernels are supported on at most a countable set and
the range is still dense.  Unitary charge transports cannot change kernels
of adjoints or closure of ranges.

\begin{theorem}[Absolutely continuous completion no-go]
\label{thm:ac-no-go}
The cofinal Hilbert completion generated from the Cauchy coefficient
module, shifted prime-power generators, Gamma graphs, covariant charge
mixing and isometric direct limits has absolutely continuous normalized
scale spectrum and no point spectrum.  Consequently every equivariant
linear map whose range vectors belong to this Hilbert target,
\[
 D:H^1_{\rm CCM}\longrightarrow\mathscr P_\infty^{\rm cov}
 \tag{10.33}
\]
annihilates every resonant eigenclass.  In particular it cannot be a
faithful descent and cannot supply the polarization of
Definition~\ref{def:global-polarization}.
\end{theorem}

\begin{proof}
The measure in (10.20) is equivalent to Lebesgue measure.  Countable
orthogonal sums, shifts of the generator, unitary conjugacy, invariant
closed subspaces, equivariant graph embeddings and isometric direct limits
preserve absolute continuity.  These are exactly the operations used in
(10.20)--(10.29); hence the target has no eigenvectors.

For every critical-line zero
\(\rho=1/2+i\gamma\), the CCM degree one has a nonzero class satisfying
\(U_t^{\rm CCM}u_\rho=e^{it\gamma}u_\rho\).  Equivariance of (10.33)
would make \(Du_\rho\) an eigenvector with the same character.  Since the
target has no point spectrum, \(Du_\rho=0\).  Infinitely many such zeros
exist unconditionally by Hardy's theorem~\cite{Hardy1914}.
\end{proof}

\begin{remark}[Hilbert-target scope]
Theorem~\ref{thm:ac-no-go} does not exclude an equivariant distributional
map into the dual member \(\Phi'\) of a rigged triple
\(\Phi\subset\mathscr P_\infty^{\rm cov}\subset\Phi'\).  Generalized
eigenvectors of an absolutely continuous operator naturally live in such a
dual.  The theorem excludes faithful descent into the Hilbert member, which
is exactly the target of the positive metric (10.27).  A resonant nuclear
construction must therefore retain evaluation vectors and jets in
\(\Phi'\) before asking which combinations admit a positive Hilbert
completion.
\end{remark}

The theorem proves a family-level statement with an explicit scope.  It
applies to decomposable absolutely continuous coefficient completions and
to the operations listed in its proof.  It does not assert that every
conceivable Hilbert space is absolutely continuous.  Its conclusion is the
relevant dichotomy: retaining transverse charge fibers creates a spurious
continuous cokernel, while removing them leaves dense range and zero
reduced cokernel.  Neither is the discrete resonant CCM degree one.

\subsection{Six terminal formulations and their common content}
\label{subsec:six-terminal}

The development reaches the same force-bearing sign through six distinct
constructions.  Their mathematical forms are different, but each becomes
a proof of Weil positivity precisely when its final gate is supplied:
\begin{center}
\renewcommand{\arraystretch}{1.2}
\begin{tabular}{>{\raggedright\arraybackslash}p{0.27\textwidth}
                >{\raggedright\arraybackslash}p{0.63\textwidth}}
\toprule
Realization & Terminal statement\\
\midrule
Localized signed forms & Lower semicontinuity of the indefinite Weil form
under the global exhaustion, not merely Mosco convergence of its positive
shift.\\
Finite inertia/Schur analysis & The physical prime--Gamma gain exceeds the
anti-shorted deficit, \(\mathfrak g_{M,J,X}>\delta_{M,J,X}\), on a
form-core exhaustion.\\
Fixed Rosati realization & A faithful factorization
\(\tau(f*g^\sharp)=\langle D[f],D[g]\rangle\); by GNS, its existence is
equivalent to positivity of the Rosati/Weil form.\\
Alternative polarization & A source-defined positive Hilbert majorant for
which normalized scaling is unitary and \(\Omega\) is bounded and weakly
nondegenerate.\\
Charged boundary pushout & Closure-faithful equivariant descent from the
complete Tate--Gamma--polar source to the separated CCM cokernel.\\
Geometric degree one & A positive global polarization
\((H^1_{\rm CCM},\Omega,J,g)\) as in Definition~\ref{def:global-polarization}.\\
\bottomrule
\end{tabular}
\end{center}

The first two rows localize the sign in finite matrices.  The third states
the sign as a factorization.  The fourth and fifth try to construct a new
positive metric before comparison with Rosati.  The last states the desired
geometric structure directly.  None of the implications from its terminal
statement to RH is missing; the missing assertion in each row is the same
global sign in the corresponding category.

This repeated endpoint admits a general formulation.

\begin{proposition}[Deformation-blindness obstruction]
\label{prop:deformation-blindness}
Let \(\mathcal X\) be a class of divisor-bearing objects, let \(\sim\) be a
deformation relation on \(\mathcal X\), and let
\(\kappa:\mathcal X\to\mathbb Z\) be an index depending on the position of
the divisor.  Suppose that a functor \(F:\mathcal X\to\mathcal Y\) is
constant on \(\sim\)-classes.  If \(\kappa\) is not constant on one such
class, then no map \(\phi:F(\mathcal X)\to\mathbb Z\) satisfies
\(\kappa=\phi\circ F\).
\end{proposition}

\begin{proof}
Choose \(x\sim y\) with \(\kappa(x)\ne\kappa(y)\).  Then
\(F(x)=F(y)\), whereas a factorization would give
\(\kappa(x)=\phi(F(x))=\phi(F(y))=\kappa(y)\), a contradiction.
\end{proof}

In the present program, Proposition~\ref{prop:deformation-blindness}
organizes four otherwise separate obstructions.  The table states the
forgetful functor actually used, the relation through which it factors,
and an explicit witness showing that the required index does not factor
through it.

\begin{center}
\small
\renewcommand{\arraystretch}{1.25}
\begin{tabular}{>{\raggedright\arraybackslash}p{0.20\textwidth}
                >{\raggedright\arraybackslash}p{0.26\textwidth}
                >{\raggedright\arraybackslash}p{0.44\textwidth}}
\toprule
Functor \(F\) & Relation \(R\) forgotten by \(F\) & Witness with different
index information\\
\midrule
Half-plane sign class of the Maass--Selberg Cauchy pairing
& Motion of resonance parameters inside \(0<\Re\rho<1\)
& Writing \(w_\rho=(1-\rho)/2\), the kernel is always
\(-1/(w_\rho+\overline{w_\sigma})\), hence negative definite.  Five
on-line resonances give signature \((0,5)\); adjoining an off-line quartet
gives \((0,9)\), still definite, although the Weil off-line index changes.\\
Scalar determinant section
& Equality of determinants after forgetting the sourced Hermitian kernel
and its Krein--Langer divisor
& The Hermitian matrices \(I_2\) and \(-I_2\) both have determinant one,
but negative indices zero and two.  The same ambiguity is why a
Nevanlinna determinant must be augmented by its sourced second variation
and divisor grading.\\
Cyclic Chern character
& Bounded operator homotopy, including scalar shifts of a finite
semilocal Weil operator
& The path from an unshifted finite Weil matrix to its lower-edge-shifted
nonnegative matrix preserves the character, while the lower edge and the
negative eigenvalue count can change.\\
Curvature of a positive determinant-line metric
& Motion of a section divisor with the Hermitian line and its metric fixed
& On a disk, \(s_a(z)=z-a\) and \(s_b(z)=z-b\) have different divisors,
but the same curvature \(\partial\bar\partial\phi\) for
\(\|1\|^2=e^{-\phi}\).  Poincare--Lelong records the moved divisor but
metric positivity does not localize it.\\
\bottomrule
\end{tabular}
\end{center}

Theorem~\ref{thm:ac-no-go} is deliberately not a fifth row of this table.
Its proof does not require a deformation witness.  It is a spectral-type
obstruction: the source has point-spectrum eigenclasses, the Hilbert target
has purely absolutely continuous spectrum, and equivariance forces every
such class to map to zero.

The conclusion is not that every analytic limit is signature-blind.
Spectral flow, for example, can retain an index when a gap and a continuous
Fredholm path are supplied.  The conclusion is narrower and operational:
none of the four functors in the table can recover the off-line index
because it factors through an equivalence relation under which that index
can change; and the absolutely continuous Hilbert completion fails for the
independent point-versus-continuous spectral reason.

\subsection{Two open programs}
\label{subsec:open-programs}

The preceding construction and no-go leave two mathematically distinct
directions.  Both must be divisor-sensitive from the start; neither may
define its positive metric from the zeros or from the Weil form whose sign
is sought.

\subsubsection{A geometric intersection product and Hodge-index theorem}

The first direction is to construct the missing geometric source of
positivity.  The finite identity (10.7) and the archimedean identity
(10.12) specify its local intersection numbers and determinant fibers.
What is still absent is a global square carrying correspondences, a
relative intersection form, and a primitive Hodge-index theorem.  A
successful object must provide:
\begin{enumerate}
\item classes \(\Gamma_{p^k}\) whose local torsion and incidence recover
\(\Lambda(p^k)/\sqrt{p^k}\);
\item an archimedean fiber whose determinant is (10.12), including the
polar \(H^0/H^2\) page;
\item a Lefschetz or supertrace formula identifying the resulting
intersection distribution with Weil's explicit formula;
\item a primitive intersection inequality whose induced Rosati form on
degree one is positive;
\item compatibility with the scaling correspondences, so that normalized
scaling is unitary on the polarized degree one.
\end{enumerate}
This is the analogue of the role played by correspondences, the Jacobian
and the Hodge index theorem in Weil's proof over finite fields.  The work of
Connes and Consani on the arithmetic curve, its Jacobian and absolute
geometry supplies the natural ambient program
\cite{CCArch,CCJacobian,CCAbsolute}.  Theorems
\ref{thm:root-torsion} and \ref{thm:gamma-polar-det} give concrete local
constraints that any such global theory must satisfy.  They do not supply
the global intersection inequality.

\subsubsection{A divisor-sensitive resonant completion}

The second direction remains analytic or derived.  It must replace the
ordinary direct-integral completion by a topology in which evaluation and
multiplicity jets at a discrete divisor survive before Hilbert reduction.
At minimum one needs a source-defined space \(\mathscr E\), a closed
restriction complex and a map
\[
 H^1_{\rm CCM}\longrightarrow H^1(\mathscr E)
 \tag{10.34}
\]
with the following properties:
\begin{enumerate}
\item (10.34) is faithful on every resonant jet and is equivariant for
scaling;
\item the Tate, Gamma and polar boundary maps act continuously and retain
the exact finite part (10.19);
\item the alternating residue form extends boundedly and nondegenerately;
\item a positive compatible complex structure is constructed from source
data, not from the zero divisor or the positive part of the Weil form;
\item the resulting normalized scale action is unitary.
\end{enumerate}
Possible ambient categories include nuclear analytic spaces, derived
cyclic modules and reproducing-kernel completions.  The category alone is
not a solution: the classical de Branges program shows that bounded point
evaluation can simply re-express the missing Hermite--Biehler sign, and
specific proposed positivity conditions for zeta and \(L\)-functions are
known to fail~\cite{deBranges1968,ConreyLi2000}.  The acceptance test is
therefore exact: the construction must fail for a Davenport--Heilbronn
type function without an Euler product~\cite{DavenportHeilbronn1936}, must preserve the CCM divisor, and
must obtain positivity from arithmetic source structure rather than assume
it.

The two programs are related but not identical.  A geometric intersection
theory could produce the resonant completion as its analytic realization.
Conversely, a successful divisor-sensitive completion might reveal the
one-dimensional shadow of the missing intersection theory.  At present
both positivity theorems remain open.

\section{Conclusion}

The Connes proof of RH, in its current mathematical form, has the following
shape:
\[
\boxed{
\begin{gathered}
\text{finite semilocal Weil self-adjointness}\\
+\ \text{finite critical-line prolate zeta-cycles}\\
+\ \text{Suzuki localized operator realization}\\
+\ \text{a faithful positive global polarization}\\
\Longrightarrow \QW\ge0
\Longrightarrow \RH .
\end{gathered}}
\tag{11.1}
\]
The first three blocks are supplied by the published
CCM/Connes/Suzuki programs in the precise sense stated in
\S\S2--7.  The finite analysis of \S\S8--9 proves that approximation,
tail control, radical shorting, prime sampling, adaptive gain recovery and
form-core exhaustion create no additional sign loss.  In that realization
the terminal assertion is the physical surplus
\(\mathfrak g_{M,J,X}>\delta_{M,J,X}\) on the heat or hybrid rows.

Section~\ref{sec:polarization-audit} then reconstructs the source geometry
behind the same assertion.  At every finite prime, a positive cyclic-root
Laplacian produces the relative torsion \(\log p\) and the incidence
\(p^{-k/2}\), hence the exact coefficient
\(\Lambda(p^k)/\sqrt{p^k}\).  At infinity, the positive Gamma spin and the
polar plane reproduce the completed archimedean determinant.  The common
boundary has an exact matched finite part, and the charge translations
assemble all prime powers equivariantly without removing cross phases.
These statements are unconditional and source-defined.

They do not complete (11.1).  Theorem~\ref{thm:ac-no-go} proves that the
ordinary absolutely continuous Hilbert completion of this source object
has no point spectrum and therefore cannot receive the resonant CCM degree
one faithfully.  Thus the finite and local polarizations are valid, while
the attempted global descent is not.  This separates a successful local
construction from the still missing global Hodge/Rosati positivity.

The six terminal formulations in \S\ref{subsec:six-terminal} are not six independent missing
lemmas.  They are six categorical expressions of the complete Weil sign.
Proposition~\ref{prop:deformation-blindness} identifies their recurring
deformation failure: scalar determinants, homotopy classes and curvature
factor through relations which forget the position or signature of the
resonant divisor.  The absolutely continuous Hilbert pushout fails for the
independent spectral reason of Theorem~\ref{thm:ac-no-go}.  Together these
results explain why further repackaging by the same operations cannot
create the missing polarization.

Two routes remain open.  The geometric route asks for an arithmetic
intersection product and a primitive Hodge-index theorem whose local
fibers are constrained by Theorems~\ref{thm:root-torsion} and
\ref{thm:gamma-polar-det}.  The divisor-sensitive route asks for a
resonant nuclear or derived completion which retains every evaluation jet
before Hilbert reduction and constructs a positive scale-compatible metric
from arithmetic source data.  Either successful construction would supply
the fourth block of (11.1).  Neither construction, nor RH, is proved here.

\begin{thebibliography}{99}

\bibitem{Weil1952}
A.~Weil,
\emph{Sur les ``formules explicites'' de la theorie des nombres premiers},
Comm. Sem. Math. Univ. Lund (1952), 252--265.

\bibitem{Connes1999}
A.~Connes,
\emph{Trace formula in noncommutative geometry and the zeros of the
Riemann zeta function},
Selecta Math. (N.S.) \textbf{5} (1999), 29--106.

\bibitem{CCArch}
A.~Connes and C.~Consani,
\emph{Weil positivity and trace formula, the archimedean place},
arXiv:2006.13771 (2020).

\bibitem{CvS2025}
A.~Connes and W.~D. van Suijlekom,
\emph{Quadratic forms, real zeros and echoes of the spectral action},
Comm. Math. Phys. \textbf{406} (2025), no.~12, Paper No.~312.

\bibitem{Connes2026}
A.~Connes,
\emph{The Riemann Hypothesis: Past, Present and a Letter Through Time},
arXiv:2602.04022 (2026).

\bibitem{CCM2025}
A.~Connes, C.~Consani, and H.~Moscovici,
\emph{Zeta spectral triples},
arXiv:2511.22755 (2025).

\bibitem{CCJacobian}
A.~Connes and C.~Consani,
\emph{On the Jacobian of \(\overline{\operatorname{Spec}\mathbb Z}\)},
arXiv:2602.15941 (2026).

\bibitem{CCAbsolute}
A.~Connes and C.~Consani,
\emph{On the Absolute Geometry of \(\operatorname{Spec}\mathbb Z\)},
arXiv:2606.06604 (2026).

\bibitem{Suzuki2025}
M.~Suzuki,
\emph{On the Hilbert space derived from the Weil distribution},
Canadian Journal of Mathematics, FirstView (2025).

\bibitem{SuzukiDirichlet}
M.~Suzuki,
\emph{Detecting zeros of Dirichlet \(L\)-functions via the Riemann
zeta-function},
arXiv:2508.17701 (2025).

\bibitem{SuzukiScrew}
M.~Suzuki,
\emph{Weil's quadratic form via the screw function},
arXiv:2606.09096 (2026).

\bibitem{Burnol2004}
J.-F.~Burnol,
\emph{Two complete and minimal systems associated with the zeros of the
Riemann zeta function},
J. Th\'eor. Nombres Bordeaux \textbf{16} (2004), no.~1, 65--94.

\bibitem{Burnol2007}
J.-F.~Burnol,
\emph{Entrelacement de co-Poisson},
Ann. Inst. Fourier (Grenoble) \textbf{57} (2007), no.~2, 525--602.

\bibitem{Deninger1998}
C.~Deninger,
\emph{Some analogies between number theory and dynamical systems on
foliated spaces},
Doc. Math., Extra Volume ICM 1998, vol.~I, 163--186.

\bibitem{deBranges1968}
L.~de Branges,
\emph{Hilbert Spaces of Entire Functions},
Prentice--Hall, Englewood Cliffs, NJ, 1968.

\bibitem{ConreyLi2000}
J.~B.~Conrey and X.-J.~Li,
\emph{A note on some positivity conditions related to zeta and
\(L\)-functions},
Internat. Math. Res. Notices (2000), no.~18, 929--940.

\bibitem{Hardy1914}
G.~H.~Hardy,
\emph{Sur les z\'eros de la fonction \(\zeta(s)\) de Riemann},
C. R. Acad. Sci. Paris \textbf{158} (1914), 1012--1014.

\bibitem{DavenportHeilbronn1936}
H.~Davenport and H.~Heilbronn,
\emph{On the zeros of certain Dirichlet series},
J. London Math. Soc. \textbf{11} (1936), 181--185, 307--312.

\bibitem{DunsterPSWF}
T.~M. Dunster,
\emph{Asymptotics of prolate spheroidal wave functions},
arXiv:1601.00699v3 (2017).

\bibitem{Krein1947}
M.~G. Krein,
\emph{The theory of self-adjoint extensions of semi-bounded Hermitian
transformations and its applications. I},
Mat. Sb. (N.S.) \textbf{20(62)} (1947), 431--495.

\bibitem{Anderson1971}
W.~N. Anderson, Jr.,
\emph{Shorted operators},
SIAM J. Appl. Math. \textbf{20} (1971), no.~4, 520--525.

\bibitem{AndersonTrapp1975}
W.~N. Anderson, Jr. and G.~E. Trapp,
\emph{Shorted operators. II},
SIAM J. Appl. Math. \textbf{28} (1975), no.~1, 60--71.

\bibitem{Haynsworth1968}
E.~V. Haynsworth,
\emph{Determination of the inertia of a partitioned Hermitian matrix},
Linear Algebra Appl. \textbf{1} (1968), 73--81.

\bibitem{Gantmacher1959}
F.~R. Gantmacher,
\emph{The Theory of Matrices},
Chelsea, New York, 1959.

\bibitem{Loewner1934}
K.~Loewner,
\emph{\"Uber monotone Matrixfunktionen},
Math. Z. \textbf{38} (1934), 177--216.

\bibitem{Bhatia1997}
R.~Bhatia,
\emph{Matrix Analysis},
Graduate Texts in Mathematics, vol.~169, Springer, 1997.

\bibitem{Karlin1968}
S.~Karlin,
\emph{Total Positivity, Vol.~I},
Stanford University Press, 1968.

\bibitem{Ismail2005}
M.~E.~H. Ismail,
\emph{Classical and Quantum Orthogonal Polynomials in One Variable},
Encyclopedia of Mathematics and its Applications, vol.~98, Cambridge
University Press, 2005.

\bibitem{Mosco1969}
U.~Mosco,
\emph{Convergence of convex sets and of solutions of variational
inequalities},
Advances in Math. \textbf{3} (1969), 510--585.

\bibitem{KuwaeShioya2003}
K.~Kuwae and T.~Shioya,
\emph{Convergence of spectral structures: a functional analytic theory
and its applications to spectral geometry},
Comm. Anal. Geom. \textbf{11} (2003), no.~4, 599--673.

\bibitem{Post2012}
O.~Post,
\emph{Spectral Analysis on Graph-like Spaces},
Lecture Notes in Mathematics, vol.~2039, Springer, 2012.

\bibitem{SuzukiUemura2014}
K.~Suzuki and T.~Uemura,
\emph{On instability of global path properties of symmetric Dirichlet
forms under Mosco-convergence},
arXiv:1412.0725 (2014).

\end{thebibliography}

\end{document}
